
\chapter{Bureaucratic Coordination in Trash Management: Evidence from Bhopal and Lucknow}

This chapter investigates the critical role of bureaucratic capacity in the provision of public goods within rapidly urbanizing societies, focusing on the cities of Bhopal and Lucknow. It argues that the effectiveness of public service delivery in these urban contexts is determined not just by overcoming social cohesion challenges but by the institutional configuration and capacity of state structures. Through a detailed analysis of Bhopal, the chapter shows how state-level bureaucrats enhance governance by effectively managing horizontal and vertical coordination, acting as boundary spanners, and leveraging institutional knowledge. In contrast, Lucknow's fragmented approach underscores the challenges posed by weaker bureaucratic frameworks. This comparative analysis underscores the importance of robust bureaucratic structures and strategic coordination in shaping state capacity and ensuring the equitable delivery of public services in diverse urban environments.

\section{Why study trash?}

Trash embodies the classical characteristics of a public good: non-excludability and non-rivalry \citep{Samuelson1954, Bennett1979}. Once waste management services are implemented, it is impractical to exclude individuals from their benefits, and one person's use does not diminish the availability of the service to others. This inherent nature necessitates governmental oversight to ensure equitable access and the overall welfare of the community. Concurrently, trash management serves as a tangible indicator of state capacity and legitimacy \citep{Fukuyama2013}. Though often perceived as a mundane aspect of daily life, trash collection reveals much about how governments operate and manage public services, reflecting the state's ability to implement policies, maintain public health, and manage resources effectively.

By efficiently organizing and delivering these essential services, governments fulfill their fundamental role in upholding public health and environmental standards while reinforcing the social contract with citizens. The visibility and universality of waste management makes it a critical area for public accountability, where performance directly impacts public satisfaction and, by extension, political stability \citep{Peters2016}.

Moreover, the study of trash collection highlights the critical role played by the ``hidden state''— local governments—in developing state capacity and providing essential public services. While often overlooked in favor of federal agencies and national policies, local governments are where the state truly meets the street \citep{Zacka2017}. Focusing on the local state allows for a deeper appreciation of the importance of local political dynamics, public-private partnerships, and the agency of actors at the margins in shaping political development. This bottom-up perspective reveals the complex and often contradictory processes through which state capacity is built and exercised.

In the case of India, described as ``weak-strong'' \citep{Rudolph1987}, ``failed developmental'' \citep{Herring1999}, ``Sisyphean'' \citep{Kapur2007} and ``flailing'' \citep{Pritchett2009} the study of trash management provides a case to untangle the mechanisms behind these characterizations. Bhopal and Lucknow are presented as ``critical cases'' \citep{Eckstein1975} for examining urban state capacity. A critical case permits analytic generalization: if a theory proves effective in the challenging conditions of a critical case, it is likely to succeed in less challenging or more typical conditions as well. These two cities, situated in the so-called 'BIMARU' states—characterized by low Human Development Index (HDI) and poor state capacity—represent a rigorous test for theories of state efficacy and institutional performance.

There is also the critical question of how much power an Indian city actually wields. Indian cities operate with a severely limited scope of authority, primarily managing tasks such as trash collection, issuing birth and death certificates, and maintaining burial grounds (Annexure-1). More significant functions, such as land use and planning, often remain under the control of para-state bodies, while even basic municipal functions like sanitation, roads, street lighting, and sewerage are only selectively delegated to city governments or shared among various agencies, creating a complex web of line bureaucracies with their own chains of command. In this context, studying trash collection and management becomes an essential lens for understanding the true functioning and capacity of the urban local state in India. It is also important to note that, by law, bureaucrats in Indian cities often possess more power than elected officials. This structural feature makes urban centers a particularly relevant focus for examining bureaucratic state capacity. If there is a case to be made for studying bureaucratic state capacity, it must begin with the cities.

\section{Descriptive Statistics and Data}

\subsection{Demographic and Administrative Data}

Table \ref{tab:admin_data} provides a comparative overview of demographic and socio-economic indicators for the two cities. Table \ref{tab:admin_structure} compares the governance frameworks of the two cities. It shows that both cities have directly elected mayors and indirectly elected city councils, with commissioners serving as the administrative heads. However, Bhopal has more autonomy, with six powers under the independent control of the city government, compared to just one in Lucknow.

\begin{table}[ht]
\centering
\caption{Administrative Data Comparison: Bhopal and Lucknow}
\label{tab:admin_data}
\begin{tabular}{|l|c|c|c|c|c|c|c|c|}
\hline
\textbf{City} & \textbf{Area} & \textbf{Population} & \textbf{Slum} & \textbf{Literacy} & \textbf{SC} & \textbf{OBC} & \textbf{Forward} & \textbf{Muslim} \\
 & \textbf{(sq. km.)} & \textbf{(2011)} & \textbf{\%} & \textbf{\%} & \textbf{\%} & \textbf{\%} & \textbf{\%} & \textbf{\%} \\
\hline
Bhopal & 286 & 1798218 & 26.68 & 83 & 14.36 & 39.98 & 42.70 & 24.69 \\
\hline
Lucknow & 349 & 2817105 & 12.95 & 86 & 10.79 & 42.26 & 43.83 & 21.79 \\
\hline
\end{tabular}
\raggedright
\small Source: Census of India 2011. The OBC, Forward caste numbers are from Brown Citizen Urban Survey
\end{table}

\begin{table}[ht]
\centering
\caption{Administrative Structure of the City Governments: Bhopal and Lucknow}
\label{tab:admin_structure}
\begin{tabular}{|l|c|c|c|c|c|}
\hline
\textbf{City} & \textbf{Mayor} & \textbf{City Council} & \textbf{Administrative} & \textbf{Powers under} & \textbf{Municipal} \\
 & & & \textbf{Head} & \textbf{independent control*} & \textbf{Budget} \\
\hline
Bhopal & Directly Elected & Indirectly Elected & Commissioner & 6 & \$26.00 bn \\
\hline
Lucknow & Directly Elected & Indirectly Elected & Commissioner & 1 & \$28.65 bn \\
\hline
\end{tabular}
\raggedright
\small *Solid waste management is common for the two cities.\\
\small Bhopal has solid waste management, fire services, urban poverty alleviation, urban amenities and facilities, burials and crematoriums. Cattle pounds, birth and death, bus stop.\\
\small Source: Respective State Municipal Acts and Budget Documents
\end{table}

\subsection{Descriptive Results from Survey Data}

The data used in this study comes from the Brown University - Urban Citizenship Survey, a face-to-face survey conducted between May and June 2022. The survey aimed to collect data on urban citizenship, across multiple cities. The sample includes 4,230 respondents. While a more detailed treatment of the data, including multivariate analysis, will be addressed in a subsequent chapter, this section focuses on providing specific observations related to disparities in waste management services.

Question asked: What is the frequency of trash collection at your house?

Table \ref{tab:trash_collection} shows that overall, Bhopal demonstrates a more systemized waste collection process, with 76.3\% of households receiving daily trash pickup, compared to only 48.63\% in Lucknow. This regularity highlights Bhopal's effective service. In contrast, Lucknow shows significant inconsistency, with 12.04\% of households reporting that their trash is never collected, indicating a higher number of households in Lucknow face inadequate waste management compared to Bhopal.

\begin{table}[ht]
\centering
\caption{Overall distribution of frequency of trash collection: City Level}
\label{tab:trash_collection}
\begin{tabular}{|l|c|c|}
\hline
\textbf{Frequency} & \textbf{Bhopal} & \textbf{Lucknow} \\
\hline
More than once a day & 11.71 & 16.82 \\
\hline
Once a day & 76.3 & 48.63 \\
\hline
Several times a week & 6.18 & 17.3 \\
\hline
Once a week & 2.47 & 2.32 \\
\hline
Less frequently than once a week & 0.87 & 1.66 \\
\hline
Never & 2.1 & 12.04 \\
\hline
Don't Know & 0.32 & 0.47 \\
\hline
Refused to answer & 0.05 & 0.76 \\
\hline
\end{tabular}
\end{table}

Tables \ref{tab:trash_caste_lucknow} and \ref{tab:trash_caste_bhopal} show significant disparities in trash collection across caste groups in Lucknow, while Bhopal exhibits more consistent service across all groups. In Bhopal, trash collection is uniform, with all caste groups receiving similar rates of daily pickup (around 74\% to 78\%), indicating an inclusive service that does not exclude any particular group.

However, in Lucknow, there is a pronounced disparity, especially among SC groups.\footnote{ST sample is excludable because of the small sample size.} A substantial 27.48\% of SCs and 51.92\% of STs report that their trash is never collected, highlighting severe inequities in service provision. Additionally, the SC and OBC groups in Lucknow exhibit a bimodal distribution: some households report that trash is never picked up, while others receive collection more than once a day. This trend is particularly stark among SC households, suggesting that while some benefit from frequent service, a significant number are completely neglected. This pattern points to a systemic issue in Lucknow's waste management, where marginalized communities, especially SCs and STs, face inconsistent and often inadequate service.

It could be that this pattern reflects the influence of networks; those who receive frequent service might be connected through networks established by successive SP and BSP governments, which historically provided jobs and access to these groups. This suggests that waste collection in Lucknow may be rationed, with access to services potentially dependent on these networks, leaving those without connections underserved.

\begin{table}[ht]
\centering
\caption{Distribution of frequency of trash collection by caste: Lucknow}
\label{tab:trash_caste_lucknow}
\begin{tabular}{|l|c|c|c|c|c|}
\hline
\textbf{Frequency} & \textbf{OBC} & \textbf{SC} & \textbf{ST} & \textbf{Gen} & \textbf{Overall} \\
\hline
More than once a day & 15.49 & 23.87 & 5.77 & 16.96 & 16.82 \\
\hline
Once a day & 51.14 & 31.98 & 30.77 & 51.55 & 48.63 \\
\hline
Several times a week & 15.83 & 11.26 & 3.85 & 21.06 & 17.3 \\
\hline
Once a week & 3.3 & 2.25 & 0 & 1.66 & 2.32 \\
\hline
Less frequently than once a week & 1.71 & 0.45 & 1.92 & 1.88 & 1.66 \\
\hline
Never & 11.5 & 27.48 & 51.92 & 6.54 & 12.04 \\
\hline
Don't Know & 0.34 & 0.9 & 1.92 & 0.33 & 0.47 \\
\hline
Refused to answer & 0.68 & 1.8 & 3.85 & 0 & 0.76 \\
\hline
\end{tabular}
\end{table}

\begin{table}[ht]
\centering
\caption{Distribution of frequency of trash collection by caste: Bhopal}
\label{tab:trash_caste_bhopal}
\begin{tabular}{|l|c|c|c|c|c|}
\hline
\textbf{Frequency} & \textbf{OBC} & \textbf{SC} & \textbf{ST} & \textbf{Gen} & \textbf{Overall} \\
\hline
More than once a day & 12.17 & 8.97 & 16.39 & 12.4 & 11.71 \\
\hline
Once a day & 77.09 & 74.09 & 78.69 & 75.87 & 76.3 \\
\hline
Several times a week & 5.25 & 9.97 & 0 & 6.03 & 6.18 \\
\hline
Once a week & 2.63 & 2.66 & 4.92 & 2.01 & 2.47 \\
\hline
Less frequently than once a week & 0.95 & 1.33 & 0 & 0.78 & 0.87 \\
\hline
Never & 1.31 & 2.99 & 0 & 2.79 & 2.1 \\
\hline
Don't Know & 0.48 & 0 & 0 & 0.11 & 0.32 \\
\hline
Refused to answer & 0.12 & 0 & 0 & 0 & 0.05 \\
\hline
\end{tabular}
\end{table}

The pattern observed in caste-based disparities is mirrored in the religious distribution of trash collection services, although the differences are smaller. In Bhopal, both Hindus and Muslims receive equal service, with similar rates of daily trash collection and low percentages of households reporting that their trash is never collected. This indicates that no religious group is particularly excluded from service in Bhopal.

In contrast, Lucknow shows notable disparities, though less pronounced than those seen in caste-based data. A higher percentage of Muslims (14.38\%) report that their trash is never collected compared to Hindus (10.66\%), and fewer Muslim households (8.5\%) receive trash collection more than once a day, compared to Hindus (18.63\%). This suggests that Muslim communities in Lucknow may be receiving less frequent and less reliable waste management services, possibly as a result of selective distribution practices. However, it's important to note that these religious differences in service access are not as stark as those observed across caste groups.

\begin{table}[ht]
\centering
\caption{Distribution of frequency of trash collection by Religion: Lucknow}
\label{tab:trash_religion_lucknow}
\begin{tabular}{|l|c|c|c|}
\hline
\textbf{Frequency} & \textbf{Hindu} & \textbf{Muslim} & \textbf{Overall} \\
\hline
More than once a day & 18.63 & 8.5 & 16.82 \\
\hline
Once a day & 48.69 & 51.2 & 48.63 \\
\hline
Several times a week & 17.42 & 18.74 & 17.3 \\
\hline
Once a week & 2.04 & 3.05 & 2.32 \\
\hline
Less frequently than once a week & 1.15 & 3.49 & 1.66 \\
\hline
Never & 10.66 & 14.38 & 12.04 \\
\hline
Don't Know & 0.45 & 0.44 & 0.47 \\
\hline
Refused to answer & 0.96 & 0.22 & 0.76 \\
\hline
\end{tabular}
\end{table}

\begin{table}[ht]
\centering
\caption{Distribution of frequency of trash collection by Religion: Bhopal}
\label{tab:trash_religion_bhopal}
\begin{tabular}{|l|c|c|c|}
\hline
\textbf{Frequency} & \textbf{Hindu} & \textbf{Muslim} & \textbf{Overall} \\
\hline
More than once a day & 11.87 & 10.78 & 11.71 \\
\hline
Once a day & 76.51 & 75.46 & 76.3 \\
\hline
Several times a week & 6.22 & 6.13 & 6.18 \\
\hline
Once a week & 2.83 & 1.67 & 2.47 \\
\hline
Less frequently than once a week & 0.94 & 0.74 & 0.87 \\
\hline
Never & 1.26 & 4.83 & 2.1 \\
\hline
Don't Know & 0.38 & 0.19 & 0.32 \\
\hline
Refused to answer & 0 & 0.19 & 0.05 \\
\hline
\end{tabular}
\end{table}

The data reveals contrasting waste management scenarios in Bhopal and Lucknow. While Bhopal demonstrates equitable trash collection across caste and religious groups, Lucknow shows significant disparities. In Lucknow, SC communities are most affected, with many households never receiving trash collection, while OBCs face challenges to a lesser extent. Religious disparities are also evident, with Muslim households experiencing more service lapses than Hindu households. These inequalities may be influenced by social networks and political connections, particularly given the historical dominance of the BSP (supported by Dalits) and SP (backed by Yadavs/OBCs) in securing sanitation jobs for their constituents.

\subsection{Overview of Bureaucratic Roles and Responsibilities in Urban Governance}

Table \ref{tab:bureaucratic_roles} gives an overview of bureaucratic roles and responsibilities in these cities. In this chapter, ``central bureaucrats'' refers to officers of the Indian Administrative Services (IAS), selected through a national competitive exam. These officers often serve as city Commissioners and Secretaries to State Governments, particularly in departments like Urban Development and Housing. In Lucknow, Commissioners are drawn from both central and state bureaucracies, while in Bhopal, they are exclusively from the central bureaucracy.

``State bureaucrats'' are officers selected through state-level exams conducted by the State Public Service Commission. These officers hold positions such as Assistant, Joint, or Additional Commissioners in city governments, and may be promoted to Commissioner roles in smaller towns and cities. Bhopal's bureaucracy includes officers from the Madhya Pradesh Municipal Administration Services, a cadre dedicated exclusively to urban administration. These career bureaucrats are posted solely in urban areas. Within the city government, they operate under the authority of the City Commissioner.

``Mid-level bureaucrats'' are technical experts in fields like engineering, sanitation, and revenue. They are either appointed by city governments or through relevant state departments, ensuring the efficient operation of city services. ``Street-level bureaucrats'' manage daily administrative tasks, including processing paperwork, maintaining records, and staffing ward and zone offices. Like mid-level bureaucrats, they are appointed either by the city governments or through relevant departments. Finally, street-level workers perform essential on-the-ground tasks such as street sweeping, trash collection, and drain cleaning. They form the largest segment of the workforce and are recruited by the city or the state urban department.

\begin{table}[ht]
\centering
\caption{Overview of Bureaucratic Roles and Responsibilities in Urban Governance}
\label{tab:bureaucratic_roles}
\begin{tabular}{|p{3cm}|p{5cm}|p{3cm}|p{5cm}|}
\hline
\textbf{Category} & \textbf{Description} & \textbf{Selection Process} & \textbf{Roles/Responsibilities} \\
\hline
Central Bureaucrats (Elite Bureaucrats) & Officers from the Indian Administrative Services (IAS), selected through a national competitive exam. & Nationally competitive exam (IAS) & Serve as city Commissioners, Secretaries to State Government departments (Urban Development, Housing, ULBs, etc.). \\
\hline
State Bureaucrats & Officers selected through state-level exams conducted by the State Public Service Commission. & State Public Service Commission exam & Hold positions such as Assistant, Joint, or Additional Commissioners in city governments, potential promotion to Commissioner. \\
\hline
Mid-Level Bureaucrats & Technical experts in fields such as engineering, sanitation, and revenue. & Appointed by city governments or relevant state departments & Provide technical expertise for efficient city operations. \\
\hline
Street-Level Bureaucrats & Handle day-to-day administrative tasks such as processing paperwork, maintaining records, and staffing ward and zone offices. & Appointed by city governments or relevant state departments & Manage daily administrative functions within city governments. \\
\hline
Street-Level Workers & Employees responsible for essential tasks like street sweeping, trash collection, drain cleaning, and operating waste management machinery. & Appointed by city governments or relevant state departments & Perform on-the-ground tasks essential for city cleanliness and public health. \\
\hline
\end{tabular}
\end{table}

\begin{figure}[ht]
\centering
\caption{Map of Lucknow City by Wards – Field Sites Highlighted in Color}
\label{fig:lucknow_map}
% Insert the figure here
\end{figure}

\begin{table}[ht]
\centering
\caption{Field Sites in Lucknow}
\label{tab:field_sites_lucknow}
\begin{tabular}{|p{2.5cm}|p{5cm}|p{6cm}|}
\hline
\textbf{Ward} & \textbf{Trash Collection} & \textbf{Rationale} \\
\hline
Aliganj & ``Contractor'' and city empaneled private concessionaire & Older Development - Salaried Class, Middle Class, Planned Colony \\
\hline
Gomti Nagar & City empaneled private Concessionaire & New Development - Upper Middle Class, Commercial establishments \\
\hline
Hazratganj & ``Contractor'' & Old Development - Significant Muslim population \\
\hline
Transport Nagar & ``Contractor'' and city empaneled private concessionaire & New Development - Migrant population, Dalits and Muslims from Bihar and Eastern Uttar Pradesh \\
\hline
\end{tabular}
\end{table}

\section{Lucknow}

In Lucknow, the elected local politicians, known as councilors, possess limited policy making power, resulting in a role that is primarily administrative rather than policy focused. Operating in a politically fragmented society where various social groups act as core patrons of political parties these politicians find themselves with little choice but to distribute patronage due to their inability to influence policy. Central politicians, who control the career incentives of Lucknow's bureaucracy, exert significant pressure on bureaucrats to comply with specific demands, aligning bureaucratic actions with political interests rather than the public good.

Consequently, trash management, being both a major employer and responsibility of the councilors, becomes a focal point for local political actors to assert their legitimacy and satisfy their patrons. Given limited city resources and the pressures from their respective patron groups, both local and elite politicians influence bureaucrats, compelling them to comply with their demands. This leads to the marginalization of the bureaucracy's capacity to pursue systemic reforms. Central bureaucrats cannot monitor their subordinates and consequently, align their career incentives with elite political leaders.

When faced with potential sanctions or reforms, lower and mid-level bureaucrats reach beyond their immediate superiors, leveraging external political connections and interdepartmental alliances. As a result, the local state chooses to privatize trash management, but that too is met with challenges. The private concessionaires are caught in the middle, navigating between local politicians, who prioritize patronage over efficient service delivery, and a constrained bureaucracy that lacks the capacity for coordination. This situation results in a failure to provide effective services.

Consequently, trash management in Lucknow highlights a scenario where fragmented societal structures shape the politician-bureaucrat relationship, ensuring that no group is entirely excluded from services. Nevertheless, the quality of these services remains substandard, impeding the development of a cohesive, city-wide system.

\subsection{Lucknow - Politics and Administration}

\subsubsection{The State}

Uttar Pradesh (UP) has been the epicenter of ``Mandal politics'' – a socio-political movement in India focused on advocating for caste-based reservations and representation.\footnote{The term ``Mandal politics'' originates from the Mandal Commission, established in 1979, which recommended reserving a significant percentage of government jobs and educational institution seats for OBCs. This movement has been central to debates on social justice, caste representation, and the political mobilization of historically disadvantaged groups in India.} This gave rise to two regional powerhouses: the Bahujan Samaj Party (BSP), championing Dalit rights, and the Samajwadi Party (SP), representing OBCs and Muslims \citep{Pai2013}.\footnote{While parties in Uttar Pradesh, as in most First-Past-the-Post systems, must build broad-based coalitions to secure electoral victories, their core support and political identity remain strongly tied to specific social groups.} These regional forces vie for political control alongside national parties like the Bharatiya Janata Party (BJP) and the Indian National Congress who's political core are the residual caste groups.\footnote{The BJP builds alliances among the non-Yadav OBC and forward castes, particularly the Rajputs whereas Congress has been traditionally seen as the party of the Brahmins.} In the 65-year history of the state, only in the last two decades have the governments completed full terms.\footnote{In the history of the state, only three Chief Ministers have completed full terms, all within this recent period. Mayawati of the Bahujan Samaj Party (BSP) served from 2007 to 2012, followed by Akhilesh Yadav of the Samajwadi Party (SP) from 2012 to 2017, and currently, Yogi Adityanath of the BJP since 2017.} The city also is a ``VIP'' seat having elected the three-term Prime Minister Atal Bihari Vajpayee (Member of Parliament from 1991 to 2009) and the current Defense Minister, a former chief of the BJP (Member of Parliament from 2014-now).

With only 22\% urbanization rate, UP has the lowest urban populations (and the highest rural population) in the country.\footnote{Census of India. (2011). Census of India 2011: Final population totals. Office of the Registrar General \& Census Commissioner, India.} This rural dominance is reflected in political representation: only 4 out of 80 parliamentary constituencies and 40 out of 403 state assembly constituencies are urban.\footnote{Computed from the Global Human Settlement layer, 2022 release. Shapefiles for the parliamentary and assembly constituencies are collected from Datameet. I considered urban centers, dense urban areas, and semi dense urban areas within a constituency as a proxy of urbanization. For methodology to formulate these categories, please see https://human-settlement.emergency.copernicus.eu/documents/GHSL\_Data\_Package\_2023.pdf?t=1716923464.} Consequently, state politicians often extract rents from cities to support activities in rural areas, explaining the underinvestment in urban areas compared to rural ones \citep{Varshney1995}.

The Dalit-supported BSP invested heavily in constructing parks across the state, including Lucknow \citep{Heath2012}. The SP focused on marquee projects like riverfront development, eight lane expressways and metro systems\footnote{Infrastructure to social development: Tracing shift in Samajwadi Party.'' Hindustan Times. March 11, 2024.} and the current BJP government is seen to invest in creating temple corridors.\footnote{Even though the Clean India campaign got a big push under the BJP, UP cities have consistently ranked the lowest in cleanliness rankings.} Focus on the fundamental infrastructure - water supply, public transport, waste management has received low priority. Lucknow's bargaining power is minimal, relegating its priorities to the bottom of the state's agenda. A IAS office in the state Secretariat candidly admitted:

``I am not ashamed to say this: everyone has used the municipality as a toolkit. Municipal ka shoshan pura administration karta hai (Municipal administration is exploited by the entire administrative system).''\footnote{Interview with Arvind Kumar Giri, Under Secretary, Uttar Pradesh. Lucknow 2022}

Administratively, Lucknow's governance structure is different from most Indian cities. The position of Commissioner in the Lucknow City Government (LMC) has been occupied by both Indian Administrative Service (IAS) officers and Provincial Civil Service (UP-PCS) officers. PCS officers are subordinate to the IAS officers in the administrative hierarchy. In contrast to most Indian cities, especially state capitals, where IAS officers typically serve as Commissioners, Lucknow only saw PCS officers in this role (until most recently). These PCS officers were more beholden to state interests, as evidenced by recent appointment patterns in Lucknow. From 1999 to 2024, Lucknow has seen 29 city commissioners. Of these 29, only five commissioners have had tenure longer than two years while 24 commissioners had an average tenure of just over six months. The tenures of the longest serving commissioners were co-terminus with the government which appointed them (Table \ref{tab:commissioner_tenures}).

\begin{table}[ht]
\centering
\caption{Recent Tenures of Municipal Commissioners of Lucknow}
\label{tab:commissioner_tenures}
\begin{tabular}{|p{4cm}|c|p{5cm}|c|}
\hline
\textbf{Party in power} & \textbf{Tenure} & \textbf{Commissioner, Lucknow City} & \textbf{Tenure} \\
\hline
Bahujan Samajwadi Party & 2007-2012 & Mr. Shailesh Kumar Singh (PCS) & 2007-2012 \\
\hline
Samajwadi Party & 2012-2017 & Mr. Udayraj Singh (PCS) & 2013-2017 \\
\hline
\multirow{2}{*}{Bharatiya Janata Party} & \multirow{2}{*}{2017-2022} & Dr. Indramani Tripathi (PCS) & 2018-2021 \\
\cline{3-4}
 & & Ajay Kumar Dwiwedi (IAS) & 2021-2023 \\
\hline
Bharatiya Janata Party & 2022-now & Inderjeet Singh (IAS) & 2023-now \\
\hline
\end{tabular}
\raggedright
\small Source: Lucknow City Government website
\end{table}

\subsubsection{The City}

Lucknow's social group includes Muslims, upper castes, OBCs, and others. Although these groups are not numerically equal, each represents a significant portion, resulting in a pluralistic demographic landscape (see Table \ref{tab:admin_data}). Social groups closely aligned with political parties as evidenced by the results of the last city election. All Muslim councilors were elected on the SP ticket, 82\% councilors in the Dalit reserved wards belonged to the BSP, and the majority of unreserved wards were won by the BJP (Table \ref{tab:political_composition}). The large forward caste population in Lucknow partly explains the BJP's dominance in the Mayoral position, which it has consistently held since 2006. Despite the limited functional power of the LMC, the political arena has been hotly contested, evidenced by the lack of a single-party majority in city elections until 2017.

\begin{table}[ht]
\centering
\caption{Political Composition of Lucknow Municipal Corporation and Corresponding State and National Governments (2000-Present)}
\label{tab:political_composition}
\begin{tabular}{|c|c|c|c|c|}
\hline
\textbf{Tenure} & \textbf{Mayor} & \textbf{Lucknow Municipal Corporation} & \textbf{State} & \textbf{National} \\
\hline
2002-2006 & Administrator & Divided - no majority & SP/BSP & NDA/UPA \\
\hline
2006 -2012 & BJP & Divided - no majority & BSP & UPA \\
\hline
2012-2017 & BJP & Divided - no majority & SP & UPA/NDA \\
\hline
2017-2023 & BJP & BJP & BJP & NDA \\
\hline
2023 - now & BJP & BJP majority & BJP & NDA \\
\hline
\end{tabular}
\raggedright
\small NDA - BJP led National Democratic Alliance, UPA - Congress led United Progressive Alliance\\
\small Source: State Election Commission for Mayor and City Government and Election Commission of India for the State Election
\end{table}

In contrast to the state and national levels, where representation for groups such as Muslims and Scheduled Castes is minimal or absent, city governments provide a more inclusive political arena. Table \ref{tab:representation} illustrates the significant gap between local-level representation and higher political offices. While Muslims make up 26\% of Lucknow's population and hold 13\% of the wards in the Lucknow Municipal Corporation (LMC), they have no representation in the State Assembly or Parliamentary constituencies covering Lucknow city. The same is true for Dalits. This disjunction between local and higher-level political representation positions the local state as a unique site for representation, particularly for marginalized communities. A city ward— the electoral unit of the LMC, characterized by a small, concentrated population—allows for better representation, enabling even marginalized groups to exert influence disproportionate to their overall population share. The local state, therefore, becomes a critical site for the expression and negotiation of group identities and interests, significantly impacting politicians in Lucknow who seek to remain competitive and secure election victories.

\begin{table}[ht]
\centering
\caption{Representation of Muslims and Scheduled Castes in Lucknow's Population, City government, and Legislative Bodies (2018-2023)}
\label{tab:representation}
\begin{tabular}{|l|c|c|}
\hline
\textbf{Level of Governance} & \textbf{Muslims} & \textbf{Scheduled Castes} \\
\hline
Population & 26\% & 10\% \\
\hline
Lucknow city government (LMC) & 13\% & 12\% \\
\hline
MLA (Assembly constituencies which form part of Lucknow) & Nil & Nil \\
\hline
MP (Parliamentary constituencies which form part Lucknow) & Nil & Nil \\
\hline
\end{tabular}
\raggedright
\small Note: The MLAs and MPs are during the fieldwork period from 2018-2023. 7 Assembly constituencies form part of Lucknow city: Bakshi Kaa Talab, Sarojini Nagar, Lucknow West, Lucknow North, Lucknow East, Lucknow Central, Lucknow Cantt. 2 Parliamentary constituency forms part of Lucknow city - Mohanlalganj PC in part, Lucknow PC in full.\\
\small Source: State Election Election Commission, Uttar Pradesh for Lucknow city data and Election Commission of India for Parliamentary and Assembly constituency data. City equivalent ACs and PCs have been derived from Delimitation of Constituencies Paper 6 and 7. Election Commission of India. Available at: https://www.eci.gov.in/delimitation
\end{table}

Lucknow exhibits significantly higher rates of non-voting political engagement compared to Bhopal, with more than double the proportion of residents holding political party memberships (Table \ref{tab:political_engagement}). This heightened political activity transforms the city into a dynamic political arena where politicians must be vigilant in guarding their territories. In the context of weak state institutions, political affiliation becomes the primary means of negotiating access to resources, further intensifying the competition among politicians to secure and deliver tangible benefits to their constituents. The strong representation of the city's diverse social groups, coupled with elevated levels of non-voting participation, means that politicians are compelled to deliver resources and services to remain competitive. The absence of formal deliberative forums in the city exacerbates this dynamic, as politicians, pressured to meet the demands of their politically active base, exert influence over bureaucrats to secure the delivery of public goods. This need to respond quickly and effectively to their constituencies' demands undermine more systematic and long-term approaches to governance, including better trash management.

\begin{table}[ht]
\centering
\caption{Political Engagement in Lucknow and Bhopal by Social Group}
\label{tab:political_engagement}
\begin{tabular}{|l|l|c|c|c|c|}
\hline
\textbf{City} & \textbf{Group} & \textbf{Time Contributed} & \textbf{Participated} & \textbf{Discussed Supporting} & \textbf{Political Party} \\
 & & \textbf{to Campaigns (\%)} & \textbf{in Rallies (\%)} & \textbf{a Candidate (\%)} & \textbf{Member (\%)} \\
\hline
\multirow{5}{*}{Lucknow} & City & 23.51 & 22.32 & 23.51 & 17.39 \\
\cline{2-6}
 & Muslims & 20.48 & 19.17 & 24.18 & 15.03 \\
\cline{2-6}
 & Hindus & 25.34 & 24.12 & 24.12 & 18.89 \\
\cline{2-6}
 & SC & 19.82 & 18.47 & 16.67 & 12.16 \\
\cline{2-6}
 & OBC & 26.08 & 25.28 & 26.08 & 19.36 \\
\cline{2-6}
 & Forward & 3.61 & 22.28 & 24.72 & 18.40 \\
\hline
\multirow{6}{*}{Bhopal} & City & 12.81 & 8.78 & 10.02 & 6.95 \\
\cline{2-6}
 & Muslims & 17.29 & 11.34 & 13.01 & 9.11 \\
\cline{2-6}
 & Hindus & 11.43 & 7.98 & 9.11 & 6.22 \\
\cline{2-6}
 & SC & 11.63 & 5.98 & 7.97 & 3.65 \\
\cline{2-6}
 & OBC & 13.60 & 9.43 & 10.38 & 5.49 \\
\cline{2-6}
 & Forward & 13.63 & 10.39 & 11.73 & 9.61 \\
\hline
\end{tabular}
\raggedright
\small The shares refer to the percentage of respondents within each social group who have said 'yes' to questions on political participation.\\
\small Source: Brown - India Citizen Urban Survey 2022
\end{table}

\subsection{Politics and Reform Resistance}

This section examines how Lucknow's fragmented political landscape influences trash management. It demonstrates that the divided social base is associated with politicians at various levels restricting bureaucratic action. The analysis highlights two key points: the persistence of patronage in waste management and resistance to reforms. These dynamics support earlier arguments that competitive representation among social groups contributes to patronage-based service delivery and reform resistance, ultimately maintaining existing power structures. The local political arena thus becomes a crucial space for identity expression and resource negotiation, often at the expense of efficient, city-wide governance.

At the street level, local councilors have maintained a tight grip on the recruitment process for crucial services such as trash collection and street sweeping. Their control over hiring practices persisted even after privatization.\footnote{Lucknow privatized its trash management in 2010 - after which it hired two concessionaires, all of whose contracts were terminated before time. More on this is discussed later.} In the early 2000s, Lucknow was no different from other Indian cities in terms of its trash management. One of the longest serving councilor of the Lucknow city government in an interview told me:

``Trash was not even something that we thought of as needing scientific solutions. In the new areas, there was some collection and disposal, but here, most of the councilors had to handle this issue. They would get their local people. It worked both ways - it gave their people employment, and got the area cleaned.''\footnote{Interview with current leader of the opposition in the Lucknow city corporation Syed Yawar Hussain Reshu, Lucknow 2022.}

In 2004, the Indian government launched the Jawaharlal Nehru National Urban Renewal Mission (JnNURM), allocating \$14 billion for urban development. This program tied funding to governance reforms, incentivizing states and cities to implement specific policies and demonstrate results to access the funds.\footnote{JNNURM established two channels for intervention: Urban Infrastructure and Governance (UIG), funding city-wide infrastructure like water supply, sewerage, urban transport, and solid waste management; and Basic Services for the Urban Poor (BSUP), targeting housing, water supply, and sanitation for the poor. Other vital sectors such as health, education, and social services fell outside its scope.} The same year, a coalition government of the BSP and the SP in the state pushed for Lucknow city government to open up hiring for sanitation workers and lower municipal bureaucracy.\footnote{Interview with Dr. P. K. Trivedi \& Dr. Shikha Singh, Giri Institute of Development Studies, Lucknow, 2022. Also read: ``When Neo-liberal Agenda Negotiates Dalit Agenda: The Case of BSP in Uttar Pradesh,'' in Trivedi, P. K. (Ed.), The Globalization Turbulence: Emerging Tensions in Indian Society (pp. 268-288). Rawat Publications.} The parties actively packed government jobs with their co-ethnics and party supporters, with the Lucknow city government being a prime target for this patronage. In Lucknow, party workers, councilors, intermediaries, and party leaders were given priority to get their ``support base'' into the LMC.\footnote{For a detailed examination of the relationship between jobs and ethnic parties, Kanchan Chandra's ``Why Ethnic Parties Succeed'' On page 85, Chandra explores how these parties leverage employment opportunities to build loyalty and support, offering a critical analysis of the political dynamics involved.} The city's bottom-heavy staff structure (5,745 out of 7,200 employees were sanitation workers) made it ideal for political patronage. While higher level positions were state-recruited, lower and street level bureaucracy along with street level workers were hired directly by the city where these jobs had minimal qualifications - just primary school education and there were no interviews.

Attempts to implement enhanced control measures over sanitation workers, including centralized attendance systems and management information systems, faced strong opposition from both bureaucrats and politicians. The first source of opposition arose from the local leaders who viewed employment in the city as one of the few avenues to secure positions for their supporters. This was particularly significant in a state where job opportunities were scarce, making government positions highly prized. Records from the LMC council meeting provide a unique view into the intersection of caste, employment, and local politics. Beni Ram Valmiki, referring to the Valmiki community, a group traditionally associated with sanitation work asked:

``Today, hiring in the Nagar Nigam is open to all castes. Traditionally, cleaning has been the occupation of my caste. How can this change be allowed? We will oppose any hiring that doesn't prioritize our community.''\footnote{Lucknow City Government meeting archives. Accessed from the Lucknow Nagar Nigam in 2021}

The assertion that certain communities should continue to perform caste-based work wasn't paradoxical; rather, it was a practical approach to ensure representation and employment in the absence of access to higher positions of power. Indeed, the administrative hierarchy even today reflects entrenched caste divisions, with upper castes predominantly occupying higher positions, while Scheduled Castes largely holding street level workers and bureaucrat's posts. Only 15\% (32 out of 212) of the officers in the LMC were Scheduled Castes (SC). While caste-wise data for sanitation workers is unavailable, it is known that the majority of the 11,200 workers come from lower caste backgrounds.\footnote{Compiled by employee names made available by the LMC. Name matching was done manually using list from the Ministry of Social Justice and Empowerment https://socialjustice.gov.in/common/76750}\footnote{Before elections, all parties make promises to Dalit groups to offer jobs in these positions. Indian Express. (2022, February 4). SP promises 40,000 jobs, hike in daily wage. Retrieved from https://indianexpress.com/article/cities/lucknow/sp-promises-40000-jobs-hike-in-daily-wage/}

In 2015, Jyoti Envirotech, a private concessionaire contracted for waste management in Lucknow, only operated in 55 out of 110 wards. The remaining wards were still managed by private contractors closely tied to local councilors and lower-level bureaucracy.\footnote{Final City Sanitation Plan for Lucknow, 2011 pg. XX} This informal division continues to date, with half the city cleaned by politician and bureaucrat-run firms, and the other half by the private concessionaire.\footnote{Hindustan Times. (2023, January 10). LMC to hire eight companies to improve Lucknow's ranking in cleanliness survey. Hindustan Times. https://www.hindustantimes.com/india-news/lmc-to-hire-eight-companies-to-improve-lucknow-s-ranking-in-cleanliness-survey-101673454369771.html} The historic Hazratganj ward exemplifies this divide, where the local councilor maintains his own sanitation team. In an interview:

``The company was only interested in managing the new area. Our area wasn't being properly serviced, so we demanded that the company withdraw and requested the Nagar Nigam to assign me 10 workers to handle the cleaning of my ward.''\footnote{Interview with Md. Saifullah, Councillor for the Yadunath Sanyal ward from 2006-2017, Lucknow 2022}

The then Additional Municipal Commissioner, PK Srivastava, acknowledged this approach, noting: ``We got the city to pay money. But given the pressure of the councilors from the old areas, we just thought of giving the contract to the local firm which we thought would do a better job.'' When pressed about whether the councilors' closeness to the then Urban Development Minister, who was Muslim, could have played a role, Srivastava smiled and cryptically responded: ``Politicians will do politics.''\footnote{Interview with PK Srivastava, Additional Municipal Commissioner LMC 2009-2013, Lucknow 2022}

The second source of opposition stemmed from the weak bureaucracy in Lucknow, forcing politicians to assume administrative roles typically handled by bureaucrats. The low capacity of the bureaucracy in Lucknow is discussed in the following section, but even quotidian tasks of filing complaints and receiving responses from the local bureaucracy, even for councilors, was difficult. In addition to not having any policy making powers (the LMC has no legislative powers) formal spaces to deliberate and sanction bureaucracy did not exist. House meetings occurred few and far between, and no mechanism for bureaucrats and politicians to work together existed.\footnote{The absence of these forums in Lucknow would make the city council meeting a display of heightened disagreements and showmanship. Of the eight meetings I attended during three years of conducting this study, only four lasted for more than an hour. In one meeting, the Commissioner threatened to kill a councilor, and in another meeting, a councilor tried to commit suicide by jumping out of the venue. While all this ruckus was occurring, no one cared about the physical infrastructure where these meetings were taking place. Amar Ujala. (2024, February 28). Ruckus in Lucknow Nagar Nigam over BJP corporator's comment. Retrieved from https://www.amarujala.com/lucknow/ruckus-in-lucknow-nagar-nigam-over-bjp-corporator-s-comment-2024-02-28} Consequently, politicians were deeply involved in the minutiae of day-to-day municipal operations. Shadowing two councilors from Gomti Nagar ward, it became apparent that their roles were primarily administrative. As one councilor candidly expressed,

``We are firefighters, trash pickers, and drain openers. That is our role (hum aag bhujane wale hai, kuda uthane wale hai, naali khulwane wale hai).''\footnote{Interview with Prithvi Gupta (BJP), councilor from Aliganj ward. Lucknow 2022}

This statement vividly illustrates how local politicians are forced to assume administrative roles typically handled by bureaucrats, directly involving themselves in mundane municipal operations. In the absence of any substantive policy-making role, councilors were compelled to rely entirely on patronage as their primary means of exercising influence and maintaining relevance. Highlighting the view of controlling sanitation workers as both a necessity for basic service provision and a display of their limited power, one councilor explained,

``Corporation never had funds to clean the city - neither then, nor now. If I do not employ my people to collect trash, my support will wane. If I don't do it, trash will not be lifted.''\footnote{Interview with Kaushal Shankar Pandey (INC) , councilor from Aliganj ward. Lucknow 2022}

\subsection{Fragmented Authority: Constraint and Complying Bureaucracy}

Politicians controlling sanitation workers might have been the handmake of elite politicians aiming to satisfy their core base, but the mechanism through which this control persists has the Lucknow bureaucrat at its center. As discussed earlier, Lucknow is unique in having Commissioners from the Uttar Pradesh Provincial Civil Services (PCS) (referred interchangeably with state bureaucracy). Unlike the IAS officer, state bureaucrat's career trajectories are confined to the state, where being appointed as the Commissioner of Lucknow, the capital city, is considered the most prized posting. In non-Commissioner postings, a state bureaucrat would report to an IAS officer, who would be their immediate senior. However, as Commissioner Lucknow, they had their own ``bureaucracy to boss over.''\footnote{Interview with MG Ramachandaran, Urban Development Secretary 2004-2009, Government of India. Delhi 2023} Consequently, bureaucratic turf wars were common, as will be detailed later in this section.

The appointment of a state bureaucrat as Lucknow's Commissioner carried significant political implications. While turnover rates remained high, the emergence of stable state governments post-2007 allowed ruling parties to appoint their preferred state bureaucrats to these positions (Table \ref{tab:commissioner_tenures}). From 1999 to 2024, Lucknow has seen 29 city commissioners. Of these 29, only five commissioners have had tenure longer than two years while 24 commissioners had an average tenure of just over six months. The tenures of the longest serving commissioners were co-terminus with the government which appointed them. Unlike IAS officers, who can opt for central postings, these state bureaucrats are confined to the state, making them more susceptible to political pressures. In Lucknow, this unique situation, dominated by state-restricted bureaucrats, compromised their ability to prioritize city needs. Instead, it heightened their responsiveness to state political elites, who leveraged these appointments to further their agendas and extract favors.

The Commissioner's reluctance to reform the sanitation worker recruitment system exemplified this political alignment. While JnNURM linked funding to reforms, including manpower accounting and standardized hiring processes, city commissioners appointed by state-level politicians, resisted these changes.

An outsider's perspective on this period in Lucknow is insightful. The federal government's strategy for overseeing reform-linked project funding relied heavily on consultants from International Financial Institutions (IFIs), including the US Agency for International Development (USAID) in Lucknow.\footnote{International Financial Institutions as part of the reform to fund projects. In Bhopal, ADB and the UK's Department for International Development (DFID) were active in urban projects. Meanwhile, in Uttar Pradesh, USAID's FIRE-D project worked on urban reforms. These IFIs promoted city development plans, market-based infrastructure financing, and capacity building reforms. In the case of Lucknow, FIRE-D USAID were responsible, and while the major funding went to water related projects, on trash.} While this approach often overlooked the entrenched political realities of the state and was obsessed with private management solutions, the core of these IFIs was to reform the public management systems in the city governments. USAID's two-year engagement with the Lucknow city government on solid waste management ended prematurely due to conflicting priorities and unmet reform expectations. USAID had set the reformation of the public management system as a prerequisite for continued involvement. In an interview, the official noted:

``Officials overseeing the project in Lucknow were continuously fighting over municipal corporation's consistent deviations from the original contract terms of forming an attendance-based system as well coming up with a double-entry accounting system''.\footnote{Interview with RK Srinivasan, USAID Water, Sanitation and Hygiene Finance (WASH-FIN) Delhi 2023}

The divergence between USAID's push for formalized human resource processes and the city Commissioner's move towards decentralized cash payments ultimately led to USAID's withdrawal.\footnote{It was in a PCS officer SP Singh's tenure as the Commissioner of Lucknow (2003-2005), that the policy was changed} The evaluation report read:

``We note very low capacity across the state and perennial recruitment challenges. Capacity to plan and execute development projects is constrained, and Lucknow's low-level staff capacity to manage solid waste management using finance and technology is limited.''\footnote{USAID. (2018). India Financial Institutions Reform and Expansion—Debt and Infrastructure Ex-Post Evaluation: WASH Ex-Post Evaluation Series—(CKM) Project. Prepared by Ecodit LLC and Social Impact Inc. under the Water CKM Project IDIQ No. AID-OAA-I-14-00069; Task Order No. AID-OAA-TO-I5-00046.}

This complying bureaucracy and informal sanitation workforce continued to have delirious effects on the city. The City Sanitation plan noted the conditions in Lucknow presciently:\footnote{Final City Sanitation Plan for Lucknow, 2011}

``Employees from the private partner collect waste in most of the parts of the city. Waste is discharged by establishments (residential and non-residential) into open plots, open drains etc. These unorganized disposal methods have resulted in the accumulation of solid waste on roadsides, vacant plots, and storm water drains. This has resulted in a number of hygiene-related problems such as breeding of mosquitoes and presence of stray animals. In the case of slum households, this role is played by a family member.''

The privatization of Lucknow's trash management system was the result of these conditions, with the contract being awarded to Jyoti Envirotech, a company owned by a prominent local contractor with established connections to the Lucknow Development Authority (LDA). The contract covered the collection, segregation, transportation, and processing of the city's waste.\footnote{Owned by Charan Jeet Singh also owned Jyoti Buildtech Private Limited.} The LDA, a state-controlled technical body headed by an IAS officer directly under the state government, wields significant influence in Lucknow's planning and development.\footnote{Including major projects such as housing, waste treatment facilities, and land allocation} Unlike the LMC, which is a political body, the LDA's technocratic structure allows it to circumvent the ``political messiness'' of the city government, facilitating closer relationships between politicians and contractors. Jyoti Envirotech's success in securing the contract was facilitated by a robust network of political and bureaucratic connections, particularly through a high-ranking bureaucrat known for close ties to the then state's chief minister.\footnote{Navneet Sehgal, a 1988 batch IAS officer was closely linked with BSP Chief Minister Mayawati, and his career progression was closely aligned with her political tenure. During Mayawati's time as chief minister, Sehgal held numerous key positions in Uttar Pradesh. His association with Mulayam Singh Yadav saw him move to central deputation, but he returned to Uttar Pradesh when Mayawati regained power. For more, Indian Express ``Navneet Sehgal: Profile of the new Prasar Bharati Chairman and former Uttar Pradesh bureaucrat.'' Indian Express. August 4, 2024. Available at: [Indian Express](https://indianexpress.com/article/political-pulse/navneet-sehgal-profile-prasar-bharati-chairman-uttar-pradesh-bureaucrat-9218146/). and India Today: ``Maya Durbar: Mayawati's core team is a mix of old and new but all share unflinching loyalty.'' India Today. August 9, 2008. Available at: [IndiaToday](https://www.indiatoday.in/magazine/cover-story/story/20080818-maya-durbar-737232-2008-08-08).} This bureaucrat's career trajectory closely mirrored the chief minister's political fortunes, at one point overseeing numerous key departments. The connection between this influential bureaucrat and the company's founder, reportedly based on shared regional heritage and prior acquaintance, was instrumental in the company's ability to secure the contract.\footnote{The professional rapport between Sehgal and Singh was strengthened by their shared Punjabi heritage and prior acquaintance from Chandigarh. Interview with Mewalal of the Muskan Jyoti Samiti. Lucknow 2023. Mewalal has been running a co-op Muskan Jyoti Samiti which provides trash management services in Lucknow from 2006. His organization was banned from participating in the tenders for trash management. His information was triangulated with retired officers of the LMC. For more on his work and how LMC cut short Mewalal's work - Profits from waste: an NGO-led initiative for solid waste management in Lucknow, Uttar Pradesh, India (English). Water and sanitation program field note Washington, D.C.: World Bank Group. http://documents.worldbank.org/curated/en/541471468041669026/Profits-from-waste-an-NGO-led-initiative-for-solid-waste-management-in-Lucknow-Uttar-Pradesh-India}

On being asked why Lucknow chose to privatize trash management, instead of building in house capacities, the then Lucknow City Commissioner, a state level PCS officer, in an interview told me:

``We were stuck from both ends. Politicians here wanted to preserve subcontracting, and pressure from Delhi and Secretariat [Urban Development Department] was to reform or give up any funding. There was no capacity with us - every drop of diesel had ``neta ka hastakshar'' (politician's signature) and ``babu ka thapa'' (a bureaucrat's stamp of approval) We thought the best way to overcome this was to privatize the system.''\footnote{Interview with SK Singh, PCS officer, Commissioner of Lucknow 2007-2012. Delhi 2023}

The privatization of Lucknow's trash management system, seen as the ``best way out'', could not overcome the deeper issues on ground. As part of the agreement, 186 waste collection vehicles were to be transferred from the city government to Jyoti Envirotech which would form the bulk of its mechanized workforce. The lower bureaucracy of LMC resisted. Lower-level bureaucrats and politicians had long exploited the management of these vehicles, particularly the allocation and use of fuel, as a lucrative source of unofficial earnings.\footnote{Dainik Bhaskar. (2024). Drivers of Municipal Corporation RR Department went on strike protesting against the arrest of a driver in the case of diesel theft. Retrieved from https://www.bhaskar.com/local/uttar-pradesh/lucknow/news/drivers-of-municipal-corporation-rr-department-went-on-strike-protesting-against-the-arrest-of-a-driver-in-the-case-of-diesel-theft-131317890.html} But more critically, no investments were made to develop auditing, regulatory, or information extraction mechanisms for the private concessionaires. This resulted in significant deficiencies in record-keeping and transparency in Lucknow's municipal solid waste management system. Without proper documentation and oversight mechanisms for private operators, it became impossible to verify whether the specified waste collection system was being implemented consistently across the city or to hold private concessionaires accountable for their contractual obligations. Furthermore, the lack of records regarding user charges, amounts due and recovered, pointed to a broader issue of financial mismanagement and revenue leakage. The audit conducted by the Auditor General of Uttar Pradesh noted:

``In NN Lucknow, records for 55 wards out of total 110, where collection of MSW was being done by the concessionaire, were not available. Thus, the audit could not verify whether the specified system was implemented in the ULBs [urban local body] for collection of waste on a regular basis… It was also noticed that NN Lucknow was not maintaining records regarding collection of user charges (details of wards, number of houses, number of bills issued, amount due and recovered as user charges). In absence of these records, the amount of due and recovered as user charges could not be ascertained by Audit.''\footnote{Office of the Principal Accountant General (General \& Social Sector Audit), Uttar Pradesh. (2016). Annual Technical Inspection Report on Urban Local Bodies for the Year Ended 31 March 2016 (p. 23, 24). Government of Uttar Pradesh. Retrieved from https://agup.nic.in/pag/en/docs/ULB\_English.pdf}

By 2015 Jyoti Envirotech faced legal issues and became a scapegoat for the city's trash management problems, leading to the termination of their contract by a court order, and the subsequent takeover by Ecogreen, a China based trash company.\footnote{A subsidiary of China Jinjiang Environment Holding Company Limited (CJE), coinciding with the Bharatiya Janata Party (BJP) government being elected.} Like Jyoti Envirotech, the change of hands to Ecogreen did not alter the ground realities. The long-standing system of parts of the city being serviced by firms associated with lower-level bureaucrats and politicians had been systematized, and my time in the field coincided with Ecogreen still servicing only four of the eight zones, similar to Jyoti Envirotech's performance.

\subsection{Internal discord, external patronage-seeking - Gomti Nagar}

The Gomti Nagar case study reveals an ecosystem where senior bureaucrats' and politicians' reform efforts are systematically constrained, paradoxically leaving those in authority powerless against their subordinates. This dynamic is rooted in a network of relationships transcending formal structures. When facing sanctions or reforms, lower and mid-level bureaucrats leverage external political connections and interdepartmental alliances, circumventing senior officials' authority. The constraint is reinforced by a system of political tutelage, where different parties protect various levels of bureaucracy. This multi-layered patronage system impedes reforms at all levels of governance, with lower-level bureaucrats protected by one political entity and state's bureaucrats by another. The result is a state structure where nearly every level enjoys some form of external political insulation against accountability measures.

Gomti Nagar: Dr. Indramani Tripathi's tenure as Municipal Commissioner coincided with the start of my fieldwork in Lucknow. Tripathi, a state level PCS officer, held office during a period of significant administrative volatility. His tenure was bookended by the rapid successions of Ajay Dwivedi and Inderjit Singh. A Brahmin, Tripathi occupied a uniquely powerful position as the President of the Uttar Pradesh Provincial Civil Service (PCS) Association - a role that carried considerable influence given that the entire state cadre was part of the association. Interviews with officials from the LMC, Lucknow Development Authority (LDA), and Urban Area Development Department (UADD) revealed that Tripathi's networks were underpinned by caste affiliations. One senior official in an interview told me:

``Dr. Tripathi's tenure was marked by palpable tensions with both state and central politicians, yet he enjoyed a degree of protection due to his relationship with Urban Development Minister A.K. Sharma.''\footnote{My fieldwork was marked by tensions surrounding Chief Minister Yogi Adityanath's perceived favoritism towards Thakurs - his co-ethnics. Media reports highlighted growing discontent among Pandits (Brahmins) who felt underrepresented in both ministerial positions and senior bureaucracy. In an apparent 'strategic' move to address these allegations, the entry of AK Sharma, a Brahmin bureaucrat, into Uttar Pradesh's politics was widely interpreted as a countermeasure. This development underscored the complex interplay of caste dynamics in the state's political and administrative spheres, reflecting the ongoing challenges in balancing representation and power distribution among different caste groups. For more, read Yadav S. The Heart of Power: The Chief Ministers of Uttar Pradesh. New Delhi: Roli Books; 2024.}

The connection between these two Tripathi's (unrelated by blood but sharing the same caste) went beyond typical bureaucratic relationships. Ram Nagina Tripathi had a checkered past - he was suspended and moved from the Ghaziabad Development Authority (GDA) due to his involvement in a ₹24 crore cable fraud case in 2015.\footnote{Times of India. (2015, June 30). Ex-GDA engineers suspended for Rs 24 crore cable fraud. Times of India. https://timesofindia.indiatimes.com/city/noida/ex-gda-engineers-suspended-for-rs-24-crore-cable-fraud/articleshow/47889426.cms} Surprisingly, this didn't hurt his career in Lucknow. Instead, he kept his job and even flourished during Indramani Tripathi's time as Commissioner.

In my field site of Gomti Nagar, I was interviewing the Executive Engineer of Zone 4 about reports of a pending National Green Tribunal inquiry into Lucknow's trash situation. The national body had wrapped up the city government and formed a committee under a former High Court judge to visit one of the legacy waste mountains in his area.\footnote{Times of India. (2022, February 18). NGT imposes Rs 10 crore fine on LMC for flouting environmental norms. Times of India. https://timesofindia.indiatimes.com/city/lucknow/ngt-imposes-rs-10-crore-fine-on-lmc-for-flouting-environmental-norms/articleshow/97998302.cms} As we were talking, there was a flurry of movement outside, and everyone started heading to the central workshop yard in the Zone office compound. Mayor Sanyukta Bhatia was conducting a surprise inspection of the central workshop (R\&R department). Surrounded by a slew of officers, the Mayor asked for fuel allocation documents. Only 11 vehicles had their entry books available, which angered the Mayor, who then demanded logbooks for all vehicles. Calling R.N. Tripathi over the phone, she asked, ``Tripathi ji, why are logbooks not being maintained? Your officer is telling me that the Nagar Nigam is taking fuel in canisters.'' As officers started piling up in the room, she continued, ``Fuel should not be provided in cans or containers under any circumstances.'' While I could not hear R.N. Tripathi's reply, the Mayor said in an angry tone, ``I will do whatever is required to maintain transparency in the Municipal Corporation,'' and cut the call.

The next day, newspapers reported on the Mayor's visit, revealing that she had requested an Economic Offences Wing (EOW) investigation into alleged corruption involving Municipal Commissioner Indramani Tripathi and Chief (R\&R) R.N. Tripathi. The letter stated, among other things,

``Even though Ecogreen is responsible for maintaining the Nigam's machines, the R\&R department is purchasing items like batteries and tires at two to four times their market rates.''\footnote{Indian Express. (2022, March 20). Lucknow mayor wants EOW to probe corruption by civic officials. Indian Express. https://indianexpress.com/article/cities/lucknow/lucknow-mayor-wants-eow-to-probe-corruption-by-civic-officials-6533614/}

Despite these allegations, both Tripathi's continued to thrive in their careers. Indramani Tripathi, after his tenure at LMC, became the head of the Lucknow Development Authority (LDA) - another organization rich in opportunities for rent-seeking. Meanwhile, R.N. Tripathi retained his position as head of the R\&R Department, seemingly unaffected by the controversy.

In conclusion, the case from Gomti Nagar illustrates how state and mid-level bureaucrats leveraged caste networks to circumvent sanctions. With a fragmented, bottom-heavy workforce controlled by councilors and a self-preserving middle and state bureaucracy, Lucknow's privatized trash management continues to face significant challenges.\footnote{Similar case existed in my second field site of Aliganj where a tax officer used his connection with the senior bureaucrats to continue working in the city government in spite of suspension order issues. I can include the case if required.}

\section{Privatization of Trash:}

Given the fragmented politics, informal workforce, and unchecked mid and top bureaucracy, the Lucknow city government opted for a wholesale transfer of responsibilities to private entities. However, the privatization process was fraught with challenges on multiple fronts. The private sector entities found themselves caught between informal employment practices and the need to navigate complex political relationships with local councilors and state bureaucrats. Consequently, the private sector's incentives aligned more with providing short-term, technocratic solutions rather than engaging deeply with the city government's internal dynamics.

For the central bureaucracy, the focus shifted towards meeting output-oriented metrics tied to federal funding programs. These included rankings in national cleanliness surveys, achieving ``Five Star'' status for waste segregation, and constructing impressive Material Recovery Facilities (MRFs). This emphasis on visible outputs rather than systemic improvements echoes the concept of ``myths and ceremony'' \citep{Mayer1977}, where organizations prioritize external legitimacy over substantive change.

The implementation of privatization led to a de facto dual system in the city, with some areas serviced by private companies and others remaining under city's control.\footnote{Office of the Principal Accountant General (General \& Social Sector Audit), Uttar Pradesh. (2016). Annual Technical Inspection Report on Urban Local Bodies for the Year Ended 31 March 2016 (p. 23, 24). Government of Uttar Pradesh. Retrieved from https://agup.nic.in/pag/en/docs/ULB\_English.pdf} This bifurcation resulted in differential wage structures for sanitation employees and inconsistent cleaning practices across the city, ultimately leading to disparate outcomes for residents. The situation became so untenable that the courts intervened, canceling both private concessionaires (Jyoti Enviro Tech and Ms. Ecogreens) contracts. Rather than learning from these setbacks, the city government responded by formally dividing the city into two operational zones. Three zones were allocated to what was officially termed ``contractors'' but were effectively politician-bureaucrat run firms, while the remaining five zones were assigned to RAMKY, a private concessionaire. This decision further entrenched the existing inefficiencies and disparities in service delivery.\footnote{Indi Newsline. (n.d.). Two ways of cleaning and trash in Lucknow. Indi Newsline. Retrieved August 6, 2024, from https://www.indinewsline.com/two-ways-of-cleaning-and-trash-in-lucknow/}

\subsection{Bureaucratic twilight zone and middle-class action}

As has been discussed, unlike most cities led by IAS officers (central bureaucrats), Lucknow's leadership has mainly been under PCS officers (state bureaucrats). This creates an unusual reporting structure. Typically, an IAS officer heading the Lucknow Municipal Corporation (LMC) would report to another IAS officer higher up. However, a PCS officer in the same position reports directly to the Secretary, bypassing their immediate IAS senior. This unconventional setup leads to conflicts and inefficiencies.

The relationship between the Lucknow Municipal Corporation (LMC) and the Lucknow Development Authority (LDA) highlights these challenges. LDA, as a separate entity, oversees land development in and around Lucknow. Since the 1990s, whenever LMC required land for waste disposal outside the city limits, it had to lease it from LDA. Additionally, any construction on this land required LDA's approval, further complicating the process due to the bureaucratic disconnect between the two organizations. While the allocation of land for trash disposal or treatment plants would typically involve straightforward coordination between the Municipal Commissioner and the District Magistrate, ideally both IAS officers operating within the same hierarchical framework, in Lucknow's case, this process was fraught with complications.

In the case of Lucknow, both private players - Jyoti and Enviro Tech's contracts, apart from door-to-door collection, required the construction of a modern municipal solid waste processing and disposal facility. The contract mandated that the existing landfill site in Shivari (a site outside Lucknow) be given to the private sector to develop this plant. However, the discord between the PCS-led LMC and the IAS-controlled LDA and Lucknow District administration led to delays in accessing land, electricity, no-pollution certificate and other permissions. This meant that Lucknow is still not able to segregate and scientifically dispose of its waste.\footnote{The National Green tribunal has formed three joint monitoring committees to monitor Lucknow's effort to solve solid waste management since 2013. The latest being Mukesh Tiwari Vs State of UP (O.A Number 08 of 2023)} This bureaucratic impasse is not merely a matter of administrative inconvenience; it has tangible implications for the city's development and the quality of life of its residents. The inability to secure land for waste management facilities meant that the city used other vacant land in its possession to build trash landfills.

The delayed allocation of the Shivari land became the linchpin in the eventual cancellation of both private concessionaires. In the case of the first concessionaire, the contract required the plant to be built within six months after signing the contract in 2010, but the land could not be obtained until 2015, turning it into a four-year ordeal that severely disrupted the project timeline.\footnote{Times of India. (2017, April 30). Mismanagement forces Lucknow to live on trash pile. https://timesofindia.indiatimes.com/city/lucknow/mismanagement-forces-lucknow-to-live-on-trash-pile/articleshow/58441019.cms?utm\_source=contentofinterest\&utm\_medium=text\&utm\_campaign=cpps} The IAS-controlled LDA and the PCS-controlled LMC became a zone of bureaucratic turf war.\footnote{But the other problem also was the weak coordination capacity of the PCS officer heading the city government. At the meetings convened for federally funded grants, the PCS cadre commissioner had the last say in matters concerning the city.} One of Ecogreen's concessionaire managers told me:

``Fundamentally the problem was getting approvals – the contract was signed between LMC and the company, but I need permission from the electricity board, pollution control board, labor commission, and land building permissions. My point persons in the city government were the Mayor and the Commissioner. For bureaucratic matters, the Commissioner's ability to prioritize and coordinate with their counterparts [meaning different line departments] just didn't happen.''\footnote{Interview with Abhishek Singh, Manager Ecogreen. Delhi 2023}

The permissions were further hampered by another set of complications: changes in state-level governments created opportunities for the new administration to potentially renegotiate kickbacks, while also giving local opposition—led by the Shivari Pradhan and backed by a local politician—ample time to mount legal challenges against the waste management facilities.\footnote{Shashikant Tiwari served as the Pradhan from 2010 to 2014. In 2023, I interviewed the current Pradhan of the area regarding the delays in the Shivari plant.} Even when permission for land came for setting up a plant about 12 kilometers away from the existing site in 2017, a dispute erupted over tipping fees between the concessionaire and the LMC, ultimately resulting in the cancellation of the concessionaires agreement.\footnote{Jyoti Envirotech sought substantial fee increases from the initial Rs 562 per metric ton set in 2011, proposing rates up to Rs 2,483 before settling on Rs 1,604.}\footnote{Times of India. (2017, April 30). Mismanagement forces Lucknow to live on trash pile. https://timesofindia.indiatimes.com/city/lucknow/mismanagement-forces-lucknow-to-live-on-trash-pile/articleshow/58441019.cms?utm\_source=contentofinterest\&utm\_medium=text\&utm\_campaign=cpps}\footnote{Writ Petition No. 35947 of 2019 (M/s Jyoti Envirotech Pvt. Ltd. vs. State of U.P. and others),}\footnote{Jyoti Envirotech sought substantial fee increases from the initial Rs 562 per metric ton set in 2011, proposing rates up to Rs 2,483 before settling on Rs 1,604.}

\subsection{Differentiated citizenship derived from differentiated services}

Delays in securing proper waste segregation land led private concessionaires and the city government to use available spaces as illegal dumping grounds. In Aliganj, a planned colony, one such site caused health issues due to unsanitary conditions and air pollution. Residents had protested since 2018 against the 150-metric ton dumping ground. When the LMC remained inactive, the Residents' Welfare Association (RWA) took legal action, moving to the High Court in 2020 and the National Green Tribunal (NGT) in 2022. The NGT imposed hefty fines on LMC and ordered the dumping site's relocation.\footnote{Times of India. (2023, April 11). Sweet scent of victory erases years of stench. https://timesofindia.indiatimes.com/city/lucknow/sweet-scent-of-victory-erases-years-of-stench/articleshowprint/99394561.cms}\footnote{National Green Tribunal, Principal Bench, New Delhi. Original Application No. 654/2022, Priyadarshini Colony D, Residence Welfare Society v. State of Uttar Pradesh \& Ors., Date of Hearing: 13.02.2023.https://greentribunal.gov.in/gen\_pdf\_test.php?filepath=L25ndF9kb2N1bWVudHMvbmd0L2Nhc2Vkb2 Nvb3JkZXJzL0RFTEhJLzIwMjMtMDItMTMvY291cnRzLzEvZGFpbHkvMTY3NjM2NzYxNjUzMzEwNzQ0MT YzZWI1NzAwN2JiMWQucGRm} In response, the city government closed the existing dumping site and moved it a few kms outside to Mohibullapur, an informal settlement where residents lacked the resources to mount a legal challenge and could only plead with the authorities for relief.

It's no surprise that the cancellation of contracts with both concessionaires stemmed from court cases filed by residents of planned colonies, who were primarily concerned about the health hazards caused by accumulating trash. However, as argued elsewhere \citep{Ghertner2015}, middle-class activism in the courts had its limitations. While Lucknow's middle class could celebrate these small victories, the more vulnerable populations, who lacked the political or legal influence, bore the brunt of the consequences. Ultimately, the city government did not address the underlying issues but merely postponed the problem.

A second result of the privatization saga was the bifurcation of the waste collection system. As has already been discussed in the previous sections, right from the enrollment of the first concessionaire, unlike the demands of the agreement for 100\% door-to-door collection, both concessionaires only covered 55 of the 110 wards in the city (5 of the 8 zones). These wards were mostly in the planned area of the city - located in the north and the east. The core unplanned area of the east and the south – where migrants and minorities of the cities lived - was operated by the politician-bureaucrat of the local state. This ``tacit'' differentiation in the city that existed in trash collection was formalized in 2023 when the tenders were invited only for a part of the city. Of the eight zones, three zones would be managed by ``contractors'', while a private concessionaire would handle the remaining five zones. As the city finalized Ramky Co. as the new concessionaire and the contractors (code word for politician-bureaucrat run firms), an officially differentiated service regime was created - labor will get different salaries, citizens in newer areas will have access to modern technology, whereas the older and poorer areas would continue to rely on manual cleaning.

In sum, the privatization of trash management in Lucknow has entrenched and exacerbated the city's existing socio-political and bureaucratic fragmentation. The division of waste collection responsibilities between private concessionaires and politically entrenched contractors has formalized the unequal distribution of public services, thereby deepening spatial and socio-economic disparities across the urban landscape. While more affluent, planned areas have seen the implementation of technologically advanced waste management systems, marginalized neighborhoods continue to rely on outdated and labor-intensive practices. This bifurcated system underscores the broader dynamics of state and urban governance in Lucknow, where political patronage and bureaucratic fragmentation subvert efforts at systemic reform. The interplay of legal challenges, political maneuvering, and bureaucratic inertia has rendered the privatization process ineffective in addressing the structural deficiencies of urban governance, instead perpetuating existing inequalities and reinforcing the socio-political stratification of the city. As such, Lucknow's experience reflects the broader challenges of urban governance in contexts marked by deep-seated political and administrative fragmentation, where efforts at privatization often serve to reproduce, rather than resolve, entrenched patterns of inequality and inefficiency.

\begin{figure}[ht]
\centering
\caption{Map of Bhopal City by Wards – Field Sites Highlighted in Color}
\label{fig:bhopal_map}
% Insert the figure here
\end{figure}

\begin{table}[ht]
\centering
\caption{Field Site in Bhopal}
\label{tab:field_sites_bhopal}
\begin{tabular}{|p{3cm}|p{4cm}|p{6cm}|}
\hline
\textbf{Ward} & \textbf{Trash Collection} & \textbf{Rationale} \\
\hline
Bhanpur & Legacy waste site & Older Development - had a 30-year-old legacy site. Comparable to Shivari site in Lucknow. \\
\hline
Maharana Pratap Nagar & Bhopal City Government & New Development - Middle and upper middle Class, Commercial establishments. Comparable to Gomti Nagar in Lucknow. \\
\hline
Bawadiya kalan & Bhopal City Government & New Development - Migrant population, Tribals and Dalits from MP, Rajasthan and Chhattisgarh. Comparable to Transport Nagar in Lucknow. \\
\hline
Teela Jamalpur & Bhopal City Government & Old Development - Significant Muslim population \\
\hline
Indira Gandhi & Bhopal City Government & Informal colony - slum. To test Bhopal City Government's capacity to cover informal settlements. \\
\hline
\end{tabular}
\end{table}

\section{Bhopal}

Bhopal not only collects and transports waste but also addresses second order problems such as waste segregation, deriving economic and energy value from waste, and managing legacy waste. Recently, Bhopal has even issued green bonds to support these efforts.

Both cities share similar political and administrative structures: they have directly elected Mayors with limited powers, Councilors without policymaking authority, and bureaucrats heading the city government. Additionally, both cities have diverse populations. However, Bhopal stands out due to unique factors within its institutional framework.

Firstly, Bhopal benefits from a specialized cadre within the Madhya Pradesh Provincial Civil Services (MP-PCS), known as the Madhya Pradesh Municipal Administrative Services (MP-MAS), which is dedicated to urban governance. These officers, allocated specifically to the municipal cadre, operate as a ``second line of command,'' utilizing their institutional knowledge and capacity to effectively coordinate within the local state. Their clear recruitment, promotion, and transfer rules enable them to build coalitions and navigate the complex relationships between politicians, elites, and lower-level bureaucrats, thereby enhancing governance and service delivery in the city. They functioned as boundary spanners \citep{Tushman1981}, connecting different sectors and facilitating cooperation.

Secondly, Bhopal's governance structure was characterized by ``islands of deliberation'' that facilitated meaningful interaction and accountability. The city council, which convened regularly in accordance with established procedural norms, provided a formalized arena where councilors could interrogate administrative actions, issue call notices, and exercise oversight over municipal processes. This institutional arrangement promoted a level of accountability within the governance framework. Moreover, zone-level meetings served as decentralized forums for targeted problem-solving, enabling bureaucrats and local politicians to engage in substantive dialogue and collaboratively address localized issues. Although these forums may not conform to the idealized notion of deliberative democracy, they nonetheless functioned as effective mechanisms for pragmatic discourse and problem-solving within the structural constraints of the system.

These ``spaces of deliberation'' and the state bureaucrats' capacity to coordinate across different government levels were key to Bhopal's success. To repeat, the state bureaucrats occupy a sweet spot in the city government's hierarchy. Their accessibility to local politicians bypassed the often-cumbersome protocols associated with central bureaucrats. Councilors could simply walk into their offices or make a phone call. Crucially, these state bureaucrats wielded enough influence to translate political concerns into administrative action. Their requests carried weight with lower-level bureaucrats, while their access to central bureaucrats — Commissioners ensured that pressing issues received high-level attention. Their ``spanning capacities - below to lower-level bureaucrats, across to their fellow cadres and local politicians and above to higher-level commissioners ensured that important issues received necessary attention, bridging the gap between political representatives and administrative machinery.

\subsection{The City}

Bhopal is the capital city of Madhya Pradesh (MP). The 1984 Union Carbide disaster left a lasting impact on its residents and development, but the city has evolved into a unique urban landscape characterized by discrete townships merged into a unified cityscape.\footnote{Union Carbide was a chemical company responsible for the Bhopal disaster in 1984, one of the world's worst industrial accidents. These include the old city, BHEL Township, Capital Project (T. T. Nagar), Bairagarh, and newer outgrowths, each with its distinct character Bhopal Development Plan 2024}

Politically, at the state level, Congress held power for decades, with the BJP dominating in recent years, interrupted briefly by a Congress term. At the city level, however, there is more competition, with both parties sharing power at different times and serving as the principal opposition (Table \ref{tab:bhopal_political}). The Congress tends to draw more support from Dalits, Tribals, and Muslims, particularly in certain city areas, while the BJP's base is stronger among OBCs \citep{Auerbach2020}. Forward castes are split between the two parties.

Bhopal's political structure is headed by a directly elected Mayor, who serves as the nominal leader. The Mayor works alongside the Mayor-in-Council (MIC), comprising representatives from the city council's majority party. The Mayor and MIC have limited authority, including the power to approve projects when at least five MIC members agree.\footnote{For a project up to \$0.6 million by the Mayor and \$1.2 million by the MIC. Only one state in India has a Mayor in Council other than MP, which is Kolkata. Kochi too has a similar structure with an empowered standing committee} This arrangement necessitates consensus between the Mayor and MIC for major project approvals. The city council consists of 85 councilors. While Bhopal, like Lucknow, lacks legislative powers, it operates under a ``Procedure for Conduct and Business Rules''.\footnote{Madhya Pradesh Municipalities (Procedure for Conduct of Business) Rules, 2005. Published vide Notification No. 25-F.1-37-05-18-3, dated 20-7-2005, M.P. Rajpatra (Asafharan) dated 20-7-2005 at pages 680 (8-15).} This framework delineates the councilors' powers, roles, and responsibilities, ensuring structured council meetings and establishing accountability mechanisms. Councilors have the right to raise questions, call notices, request information, and audit the administration.

\begin{table}[ht]
\centering
\caption{Political Composition of Bhopal Municipal Corporation and Corresponding State and National Governments (2000-Present)}
\label{tab:bhopal_political}
\begin{tabular}{|c|c|c|c|c|}
\hline
\textbf{Tenure} & \textbf{Mayor} & \textbf{Bhopal Municipal Corporation} & \textbf{State} & \textbf{National} \\
\hline
2000-2004 & Congress & Congress & BJP & NDA \\
\hline
2004-2009 & Congress & BJP & BJP & UPA \\
\hline
2009-2014 & BJP & Congress & BJP & UPA \\
\hline
2015-2020 & BJP & Congress & BJP, Congress & NDA \\
\hline
2022-present & BJP & BJP & BJP & NDA \\
\hline
\end{tabular}
\raggedright
\small NDA - BJP led National Democratic Alliance, UPA - Congress led United Progressive Alliance\\
\small Source: State Election Commission for Mayor and City Government and Election Commission of India for the State Election
\end{table}

The Bhopal City Government's (BMC) administrative head is the Commissioner, who is an Indian Administrative Service (IAS) officer.\footnote{Government of Madhya Pradesh. The M.P. Municipal Corporation Act, 1956. Indian Kanoon. 1956. Available at: https://indiankanoon.org/doc/173401616/} The average tenure across these years for municipal commissioners has been approximately 19 months (Table \ref{tab:bhopal_commissioners}). There has not been a consistent alignment of the commissioners' tenure with the full term of the political party in power, suggesting frequent changes in bureaucratic appointments. As previously discussed, Bhopal's bureaucratic structure includes officers who are part of a specialized cadre dedicated to urban administration (MP-Municipal Administration Services).\footnote{Madhya Pradesh Municipal Services Scale of Pay and Allowances Rules. (1967). Madhya Pradesh Municipal Services Scale of Pay and Allowances Rules, 1967. Retrieved from https://www.latestlaws.com/bare-acts/state-acts-rules/madhya-pradesh-state-laws/m-p-municipalities-act-1961/m-p-municipal-services-scale-of-pay-and-allowances-rules-1967} These officers are recruited through a competitive state examination and benefit from a well-defined system of promotions, transfers, and a transparent salary structure. As career bureaucrats, their postings are exclusively within the urban areas of the state, including cities, towns, and transitional urban zones.\footnote{The following State Services exist for Municipal governance i. MP Municipal Administration Service ii. MP Municipal Engineering Service iii. MP Sanitation Service iv MP Municipal Finance and Revenue Services.} Within the city government, they operate under the direct authority of the City Commissioner, ensuring a focused and consistent approach to urban governance.\footnote{Termed MP- MAS (Madhya Pradesh Municipal Administrative Service) services such as the State Urban Administrative Service, State Urban Engineering Service, State Urban Sanitation Service, State Urban Revenue Service, and State Urban Finance Service Madhya Pradesh is one of the few states in India with Odisha, Maharashtra, Andhra Pradesh, Telangana, Tamil Nadu.}

\begin{table}[ht]
\centering
\caption{Average Tenure of Commissioner Under Different Political Parties in Power (1998-2023)}
\label{tab:bhopal_commissioners}
\begin{tabular}{|l|c|c|c|}
\hline
\textbf{Party in Power} & \textbf{Tenure} & \textbf{Number} & \textbf{Average Tenure} \\
\hline
Indian National Congress & 1998-2003 & Two & 30 months \\
\hline
Bharatiya Janata Party & 2003-2008 & Three & 20 months \\
\hline
Bharatiya Janata Party & 2008-2013 & Two & 30 months \\
\hline
Bharatiya Janata Party & 2013-2018 & Three & 16 months \\
\hline
Indian National Congress & 2018-2020 & Two & 12 months \\
\hline
Bharatiya Janata Party & 2020-2023 & Four & 9 months \\
\hline
\end{tabular}
\raggedright
\small Source: Author compiled data from field work and Department of Personnel and Training. (n.d.). Executive record sheet of IAS officers. Retrieved from https://persmin.gov.in/supremohelp/KnowYourOfficerIas.htm
\end{table}

Administratively, Bhopal is divided into zones, each encompassing approximately five wards. Ward committees are established at the zonal level, consisting of 12 councilors and one chairperson. These committees convene monthly, resulting in six annual meetings.

\section{Politics}

In 2007, Bhopal's trash management infrastructure was rudimentary. The city employed 2,780 workers, of whom 1,900 were permanent employees and 500 were daily wage laborers, supported by 50 vehicles.\footnote{Bhopal City Development Plan. Jawaharlal Nehru National Urban Renewal Mission. Bhopal Municipal Corporation; Mehta \& Associates Architects \& Planners, Indore. Published under JnNURM. 2007} The city's zones operated under a decentralized waste collection system, supervised by health officers assisted by ward-level inspectors (Daroga) and sanitary supervisors. The sweeping relied on manual labor equipped with basic tools: wheelbarrows, long brooms, and panzers. Workers, assigned to specific ``beats,'' transported collected waste to overwhelmed storage sites known as collection points.\footnote{ibid.}

Concurrently, the national Congress-led government wielded influence through JnNURM by providing funds to cities in exchange for implementing reforms in city governance. At the local level, Bhopal's directly elected mayor, also from the Congress party, advocated for the effective implementation of these schemes within the city.\footnote{Bhopal: Congress manifesto promises to half property, water tax, dust-free Bhopal. Free Press Journal. June 08, 207. Available at: https://www.freepressjournal.in/bhopal/bhopal-congress-manifesto-promises-to-get-jnnurm-money-for-waste-management-bhopal}

The alignment of political interests during this period was facilitated by a notable configuration of key figures in local and state government. Sunil Sood, a prominent Congress party member and active city politician, served as Bhopal's Mayor from 2004 to 2009.\footnote{Government of Madhya Pradesh. MP State Election Commission. MP Local Election. 2024. Available at: https://mplocalelection.gov.in/} Simultaneously, Babulal Gaur of the BJP, who was an MLA from a constituency in Bhopal, held the position of Chief Minister from 2005 to 2008.\footnote{India Today. (2019, August 21). Babulal Gaur: The man who could never believe he was king. Retrieved from https://www.indiatoday.in/india/story/babulal-gaur-man-who-could-never-believe-he-was-king-1590058-2019-08-21} The subsequent Chief Minister, Shivraj Singh Chaohan, hailing from nearby Budhni, also prioritized Bhopal's interests.\footnote{Chouhan, S. S. (n.d.). Transforming Madhya Pradesh. Skoch Group. Retrieved from https://books.skoch.in/shop/books/transforming-madhya-pradesh-shivraj-singh-chouhan/} This political continuity was further reinforced when Sood's mayoral term was succeeded by Babulal Gaur's daughter-in-law. Sood's account of this period is particularly illuminating:

``Mukhyamantri dusri dal se the aur hamare aur unke beech tana tani rehti thi, parantu kam karane ke liye kai baar dhara ke vipreet chalna padta tha. Vision align karna padta tha aur negotion karne padte the'' (The Chief Minister was from another party, and there was tension between us, but sometimes, to get work done, you have to go against the current.)\footnote{Interview with Sunil Sood, Congress led Mayor of Bhopal City Government (2004-2009), Leader of Bhopal City Council (2000-2004),, City Councillor (1995-2009)}

Former Mayor's reflections of the time highlight several key dynamics that shaped the city's approach to trash management. Central to Sood's strategy was the importance of engaging politicians to navigate funding opportunities from national schemes like JnNURM and IFIs. He stressed the risk of losing these financial resources if politicians failed to find a way to collaborate effectively. Throughout multiple interviews, he consistently returned to the idea of aligning ``vision and negotiation,'' reflecting a dual approach: a normative commitment to collaborative development efforts, coupled with a pragmatic acceptance of political realities, including the allocation of ``cuts'' from rents to various political actors.

This agreement, however, was designed into the system. The MIC, the ``Conduct of Business and Rules'' for scrutiny were tools available to the opposition, serving as crucial factors in shaping decision-making processes. This structure frequently led to deadlocks, as the opposition utilized its power to obstruct schemes. However, the potential for financial repercussions, coupled with pressures from both state and central governments, often compelled local political figures to seek compromise. The balance of obstructive powers among various stakeholders required negotiation, and therefore necessitated consensus among multiple parties. The Mayor from another term observed:

``The union government and the opposition were of the same party and had the 'veto' power but so did I we had to reach a consensus.''\footnote{Interview with Krishna Gaur, BJP led Mayor of Bhopal City Government (2009-2013), MLA (2023-now).}

A significant distinction between Bhopal and Lucknow is the extent of political representation afforded to marginalized social groups. In Bhopal, Muslims and Scheduled Castes are represented at both the city government (BMC) and state assembly (MLA) levels, unlike in Lucknow, where data indicates no representation for these groups at either level. This is notable given that both cities have similar population percentages for these groups, yet Bhopal offers them more substantial political representation (Table \ref{tab:bhopal_representation}).

Bhopal's structure likely provides marginalized groups with greater avenues for political engagement beyond the municipal level. This representation creates opportunities for these communities to be visible, acknowledged, and potentially have their concerns addressed within the political process. The inclusion of these groups in decision-making processes may foster a more inclusive governance approach, leading to more effective and equitable solutions for city-level issues that impact these communities. The principal difference, therefore, is that politicians in Bhopal are not compelled to engage in the same level of intense patronage as those in Lucknow. Due to Bhopal's more inclusive political representation, particularly of marginalized groups, there is less pressure on politicians to rely heavily on patronage networks to secure their political base. This broader representation allows for a governance structure that is less dependent on patronage-driven practices, in contrast to the situation in Lucknow, where the lack of representation for certain groups necessitates a greater reliance on patronage to maintain political support.

\begin{table}[ht]
\centering
\caption{Representation of Muslims and Scheduled Castes in Bhopal's Population, City government, and Legislative Bodies (2018-2023)}
\label{tab:bhopal_representation}
\begin{tabular}{|l|c|c|}
\hline
\textbf{Level} & \textbf{Muslims} & \textbf{Schedule Caste} \\
\hline
Population & 26\% & 10\% \\
\hline
Bhopal city government (BMC) & 17.65\% & 14\% \\
\hline
MLA (Assembly constituencies which form part of Bhopal) & 33\% & 16.67\% \\
\hline
MP (Parliamentary constituencies which form part Bhopal) & Nil & Nil \\
\hline
\end{tabular}
\raggedright
\small Note: The MLAs and MPs are during the fieldwork period from 2018-2023. 6 Assembly constituencies forms part of Bhopal city: Bhopal Uttar, Narela, Bhopal Dakshin Paschim, Bhopal Madhya, Govindpura, Huzur. 1 Parliamentary constituency forms part of Bhopal city\\
\small Source: State Election Election Commission, Madhya Pradesh for Bhopal city data and Election Commission of India for Parliamentary and Assembly constituency data. City equivalent ACs and PCs have been derived from Delimitation of Constituencies Paper 6 and 7. Election Commission of India. Available at: https://www.eci.gov.in/delimitation
\end{table}

\section{``Intermediaries inside the State''}

The management of sanitation workers in Bhopal faced challenges similar to those in Lucknow. Sanitation workers, comprising the largest employee group and a significant portion of the city government's expenditure, were plagued by irregular recruitment practices.\footnote{Saharwar, J. S. (2007). Planning proposals for solid waste management of Bhopal city (Master's dissertation). Department of Architecture and Planning, Indian Institute of Technology Roorkee. Retrieved from http://shodhbhagirathi.iitr.ac.in:8081/jspui/bitstream/123456789/12740/3/G13499.pdf} Appointments were predominantly ``ad hoc'' and, in the absence of formal regulations, were often influenced by political interventions.\footnote{ibid.}

The allocation of \$10 million to Bhopal under the JnNURM for ``urban environment'' improvement, intensified pressure on the sanitation workforce.\footnote{Hindustan Times. (2006, March 30). Rs 90 crore approved in the first phase of JNNURM. https://www.hindustantimes.com/india/rs-90-crore-approved-in-first-phase-of-jnnurm/story-03yPTybB76P5Q83e5Glhk.html. Accessed July 30, 2024.} The allocation necessitated visible reforms in manpower management and outcomes, the pressure of which was evident in the observations of a sanitation supervisor:

``There was no established system for assigning tasks. Some days involved sweeping an entire ward, while other days involved no sweeping at all. I had to hire people some days and other days we were idle. So compulsory attendance would not work''\footnote{Interview with Rajesh Kuril, Assistant Health Officer in charge Zone 3, Bhopal City Government, and previously the Safai Daroga (Inspector) in the area. Bhopal 2022}

The city experienced two strikes by sanitation workers demanding laxity in attendance measures, which adversely affected service quality and drew the ire of the Mayor, who demanded additional funds for the city from the Chief Minister.\footnote{India Today. (2006, July 15). Bhopal mayor Vibha Patel barged into a meeting on the city water crisis, called on firebrand. India Today. https://www.indiatoday.in/magazine/indiascope/story/20020715-bhopal-mayor-vibha-patel-barged-into-a-meeting-on-the-city-water-crisis-called-on}

An Assistant Commissioner from the state bureaucracy highlighted the daily challenges faced by field officers.\footnote{Interview with Neelesh Dubey, Madhya Pradesh Municipal Administrative Services. Bhopal 2022} Frontline workers petitioned for workforce expansion, rationalization of cleaning schedules, and salary enhancements, regularly bringing these issues to the attention of their superiors - the lower-level bureaucrats.\footnote{Health Officers, previously known as sanitation in-charges, hold an official position in the Bhopal City Government. They are responsible for solid waste management and drainage systems.} These lower level bureaucrats in turn brought the complaints to their senior - the state level cadre repeatedly. Assistant Commissioner from the state bureaucracy, provided a broader perspective on the issue:

``This was not just a Bhopal problem – unlike other states in India, MP had many growing cities and we had this problem coming up in Gwalior, Indore, Jabalpur, Ujjain.. And we regularly attempted to bring the staffing issue to the attention of Commissioner Mam but that is not their top priority. No one cared because they didn't go to the field. Their field experience is only when they are District Magistrate (DM). But you tell me what portion of a district is even urban? Woh municipal ko samjah hi nahin sakte (They cannot understand cities).''\footnote{Interview with Neelesh Dubey, Madhya Pradesh Municipal Administrative Services. Bhopal 2022}

His statement underscores a critical issue: the Indian Administrative Service (IAS) officers often lacked a comprehensive understanding of urban areas. Their field training primarily focused on district administration, where urban areas constituted only a small portion. This rural-centric perspective persisted even when they ascended to city commissioner roles, resulting in limited interest or expertise in urban waste management and manpower issues. At the level of elite bureaucrats, urban sanitation workforce management was often overshadowed by other priorities. This is where officers of the State bureaucracy came into focus. On the other hand, the state bureaucracy had their career trajectory linked to serving only in the city. They started as Manager (Field) and moved up the ladder and at the far end of their career were made Commissioners in small or mid-sized cities reflecting a career long investment in the urban administration only.\footnote{The career trajectory began with the role of manager, followed by assistant manager. Subsequently, they advanced to the position of Assistant Commissioner, during which their grade was upgraded (from Second Grade to First Grade and then to Selection Grade). Following this, they were promoted to Deputy Commissioner, and ultimately to Municipal Commissioner.} In the course of my fieldwork, I observed frequent interactions among Madhya Pradesh urban cadre (MP-MAS) officers posted in Bhopal.\footnote{Nilesh Dubey, Devendra Singh Chauhan, Abhay Ranjan Gowankar, and Manan, a few of the officers}

Reflecting to the earlier question of how reforms in the sanitation workers recruitment took place, another Assistant Commissioner from the State bureaucracy explained:

``We saw a problem and acted on it. The Safai Daroga (cleaning inspectors) would come and complain, and the sub-engineers too started writing letters. We had to face them, not those in the Secretariat (Sachivalaya).''\footnote{Interview with Manan, Madhya Pradesh State Urban Administration Cadre Bhopal 2021}

In response to the crisis involving striking sanitation workers in Bhopal and pressure from politicians, the State bureaucracy took proactive steps to reshape service rules governing the management of sanitation workers. Their efforts centered on standardizing departments and posts across all city governments in the state, with particular emphasis on establishing population-based staffing norms and compensation structures. They also leveraged elite networks by enlisting the support of M.N. Buch, a retired IAS officer who commanded significant respect within the senior bureaucracy. This move mitigated potential sanctions from higher authorities\footnote{Interview, Nirmala Buch, IAS and former Chief Secretary of Madhya Pradesh, who is also the wife of M.N. Buch runs the National Centre for Human Settlements \& Environment (NCHSE). Bhopal 2021} A joint letter from four assistant commissioners of the state bureaucracy and M.N. Buch, addressed to the Commissioners of Bhopal and Indore City governments and the Principal Secretary of the Madhya Pradesh government, articulated a range of operational, infrastructural, and human resource concerns in waste management.\footnote{National Centre for Human Settlements and Environment. (n.d.). Projects list: GIS. Retrieved from https://www.nchse.org/projectlist-gis.html} The letter stated:

``One of the core problems of the ULBs (urban local bodies) in India has been lack of competent and appropriate functionaries. Strengthening the capacities of the municipal staff has been a long neglected area. Even though numerous steps were taken to improve the situation, a comprehensive perspective to improve the situation was missing. There was a need to bring in more uniformity in personnel structure, clarity of roles and regular skill upgradation of employees in the ULBs.''\footnote{Department of Urban Development, Madhya Pradesh. (2008). Letter number MP-UADD/734/2008. Department of Urban Development, Madhya Pradesh.}

The joint letter highlighted the arbitrary nature of trash management, which negatively impacted worker morale and operational efficiency. The document also pointed out significant infrastructural deficits, advocating for modern equipment like tipper vehicles and waste management bins. Furthermore, it brought attention to the long-neglected aspects of worker welfare, including stagnant salaries and the absence of mandated benefits. Lastly, the letter emphasized the necessity for a clear transfers and promotions policy to enable career advancement within the sanitation workforce.

The central bureaucracy's initial response to these concerns was to engage a civil society organization for Geographic Information System (GIS) mapping of Bhopal's solid waste infrastructure.\footnote{National Centre for Human Settlements and Environment. (n.d.). NCHSE GIS. Retrieved from https://www.nchse.org/nchsegis.html} However, this approach was perceived as inadequate by the state bureaucrats. As one officer noted:

``GIS map was helpful, but what we really needed was understanding why can't our workers get into the bylanes of purana sheher (old Bhopal). Why do workers employed by us subcontract to ghost employees? For us, these were more important questions.''\footnote{Interview with Manan, Madhya Pradesh State Urban Administration Cadre Bhopal 2021}

Consequently, the state bureaucrats adopted a more hands-on approach. They collaborated across two cities to conduct field-level assessments of resource requirements and logistical challenges.\footnote{Ibid. Also see Praja's report on the evolution of reform in Madhya Pradesh in general Praja Foundation. Urban Governance Study: Madhya Pradesh Consultation. Praja.org. February 5, 2019. Available at: https://www.praja.org/praja\_docs/praja\_downloads/Urban\_Governance\_Study\_Madhya\_Pradesh\_Consultation.pdf} This fieldwork informed the development of an internal document that established operating procedures and a population-to-area-to-sweeper ratio, providing an objective measure for determining staffing needs. This document included meticulous forecasts of solid waste generation and calculated requirements for sanitary workers, vehicles, and equipment.\footnote{ibid.}

Reforming the state, especially from within, is a challenging process. The case highlights the critical role played by state bureaucrats in enhancing state capacity and driving institutional change within the Indian administrative system. Despite the initial lack of tangible results, these bureaucrats demonstrated their ability to navigate complex organizational structures and advocate for reforms. They bridged the gap between frontline workers and senior leadership, translating ground-level knowledge into actionable insights. Additionally, they collaborated with their peers across cities to organize and push for improvements. While the process of reforming the state is challenging, especially from within, this case underscores the significance of state bureaucrats as key agents of change.

\subsection{Boundary Spanning Role of State Bureaucracy}

In the preceding sections, I presented two arguments. First, I argued that the local politicians had some constraining power over bureaucrats. The legal rights of a local politician to audit records, question the functioning of the bureaucracy in the city government coupled with the practical zone level meetings where the local politicians sat with the State bureaucrat allowed a degree of bargaining power to them. Second, I demonstrated that the State bureaucrat possessed a significant capacity to organize internally and coordinate within the local state. This section examines the interactions between these two. It explores what happens when reforms within the state necessitate political intervention. Using the case of trash management in Bhopal this section argues that the effort to formalize sanitation workers through the implementation of a Management Information System (MIS) proved to be intensely controversial. Local politicians exerted considerable pressure to ensure that these workers remained under their control. In this context, State bureaucrats emerged as key figures in navigating the complex terrain of reform implementation and political resistance.

Even though the city had formalized the hiring process and created workflows for 100\% door to door collection in 2011, the actual realization of these reforms was at best sketchy.

``Despite the passage of the Municipal Workers Manual and Regulation Act four years ago, Bhopal has yet to fully realize its intended benefits. The Act's provisions for streamlined waste management and worker formalization remain largely unimplemented. While a comprehensive legal framework exists on paper, its practical application on the ground has been minimal. An audit of eight study sites revealed significant issues: numerous workers were found to be absent, and waste collection occurred only once every four days.''\footnote{Municipal Corporation (2013). Municipal Corporation Act, Project: and Local Self Government, M.P., Bhopal; Documentation on Access to Municipal Services. Retrieved from https://www.nchse.org/projectslist.htm}

Seeing manpower management as the core inhibiting problem for cities to deliver better public goods, International Financial Institutions wanted the implementation of ``modern management practices'' to regulate municipal employees in Bhopal, with a particular focus on sanitation workers. These reforms were established as prerequisites for cities to receive funding in subsequent years, mirroring the situation in Lucknow. Central bureaucracy, along with the ADB, advocated for the implementation of these management solutions to enhance city management. The SAP business solution promised to ``streamline core business processes including finance, human resources, procurement, and supply chain management.''\footnote{SAP (Systems, Applications, and Products in Data Processing) refers to a suite of software solutions designed to help government organizations and public institutions manage their operations more efficiently} The implementation of these reforms was designated as the responsibility of the city government, and, as in Lucknow, there was a demand for establishing a Program Implementation Unit (PIU). However, the response to this demand differed between the two cities. While Lucknow's city government chose to forgo the grant rather than implement the PIU within their municipal structure, Bhopal took a different approach. The IAS officers agreed to house the PIU within its organizational framework, with State bureaucrats appointed to head these newly formed units.\footnote{The key driver of this reform was Deepti Gaur Mukerjee, a 1993 MP cadre officer. She served as Commissioner of Bhopal in 1997-98, Director of Town and Country Planning from 2006-2008, led the ADB-DFID Project Uday from 2008-2010, and was Director of Housing and Urban Poverty Alleviation in the Government of India from 2010-2013.}

Before diving into the details of how politicians bulwarked the implementation of these tools, a slight detour of the unique nature of IFIs supported PIU in Bhopal is required.

In Bhopal, Project (UWSEIP), acronym Project UDAY, funded by the ADB and DFID, aimed to ``improve urban infrastructure and services and strengthen the capacities of local authorities in the areas of resource mobilization and cost recovery \citep{Auerbach2020}. Operating from 2005 to 2015, Project UDAY focused on establishing, recruiting, and building capacities within the Bhopal city government. In the hierarchical Indian state, Project UDAY held significant clout, with the state's Chief Secretary as its chairperson and mandated quarterly meetings. However, the real muscle and gut of the project was the Program Implementation Unit (PIU), embedded in the Bhopal city government. It was a 'creative' bureaucratic entity. The PIU brought all line departments under one roof for ease of coordination and was headed by members of the State bureaucracy. Styled with a corporate flair, the PIU designated its officers with titles such as Project Managers, Systems Officer, and Customer Engagement officers. It also introduced systems considered ``alien to the Indian bureaucratic arrangements''\footnote{Project Uday. (n.d.). Project management. Retrieved from https://web.archive.org/web/20180828214026/http://projectuday.gov.in/ProjectMgt.htm} by bringing in procurement, management information, and a Monitoring and Evaluation officer - all recruited from within the state cadre. The PIU's mandate expanded beyond the usual remit, incorporating social and community development officers and urban governance specialists.\footnote{The Project Implementation Unit (PIU) was overseen by a City Steering Committee (CSC), which had a broad representative structure. The CSC included political figures (Mayor, Leader of Opposition), stakeholders (Pollution Board, Chamber of Commerce), bureaucrats, and non-governmental organizations (a woman involved in social work in the city). The CSC was mandated to meet at least quarterly to monitor and guide the PIU's performance, ensuring a multi-stakeholder approach to urban governance and project implementation.} The capacity building of State bureaucrats, while not explicitly stated as a core strategy of Project UDAY, proved to be a crucial component in developing a reservoir of talent within the local government structure. As we will see in the other empirical chapter on water, Project UDAY shaped bureaucratic norms within the cadre and generated valuable institutional knowledge, significantly enhancing the Bhopal local state's capacity to deliver public goods effectively. The continuity of these officers within the urban governance system, despite transfers between cities, has been a key factor in this success.\footnote{Interview with Bruce Pollock, who headed the DFID program in Madhya Pradesh from 2008 to 2013. As one observer noted, ``These officers moved from one city to another but were retained in the PIU, so we knew we could rely on training a select group of officers which, even if transferred, would remain part of the system.'' Zoom 2021, 2022.}

While the union government required the management tools incorporated in Bhopal sooner, the PIU as the ``nodal agency'' took lead. The BMC partnered with a private entity to integrate services across all 70 wards into a unified, centralized system.\footnote{Deloitte was selected from a pool of bidders.} Bhopal, along with Mumbai and Ahmedabad, was among the early cities in India to implement a transition known as ``MIS'' (Management Information System). This change led to notable alterations in municipal operations, with a particular impact on trash management systems.\footnote{People in the field colloquially refer to this as just ``MIS.''}

The introduction of the MIS in Bhopal was met with considerable resistance, primarily due to its disruptive impact on long-established practices of the councilors. This resistance was rooted in the restructuring of traditional channels for addressing citizens' concerns, which had historically centered around local councilors and zonal offices. Bhopal, like most Indian cities, operated on a highly localized system of governance and problem-solving. The councilor and the zonal office served as crucial points for all citizen-related issues ranging from clogged drains to trash collection.

The operational mechanism was straightforward yet effective: citizens would approach their councilor or the 'chota neta' (small-scale leaders in slums or Resident Welfare Associations) with their grievances. These local leaders would then leverage their connections within the Nagar Nigam (Municipal Corporation) to address the issues \citep{Auerbach2019}. This system was deeply rooted in personal relationships and local knowledge. However, the introduction of the MIS significantly altered this landscape. By moving everything online, it created a more centralized system that, while theoretically more efficient, exposed several practical challenges. Many zonal offices lacked basic infrastructure such as computers, creating significant barriers to accessing the new system.\footnote{Interview with Prakash Narayn, Infrastructure Officer, Project UDAY Bhopal 2023} This shift to an online system potentially excluded citizens who were not technologically savvy or lacked access to digital resources. It also removed the personal element from problem-solving, which had been a key feature of the previous arrangement.\footnote{ibid.} Importantly, councilors and 'chota netas' found their traditional role as problem-solvers diminishing, potentially affecting their political influence.

Expectedly, the project ran into significant challenges. The then Commissioner acknowledged:

``Each ward had its own style of working and not having an immediate uptake on the centralized system was expected but when the Mayor led the protests, I knew this was not going to work''\footnote{Interview with Vishesh Garhpale, Indian Administrative Services, Madhya Pradesh Cadre. Bhopal 2022}

The resistance to the MIS came from both ends - the treasury benches and the opposition. At the monthly Bhopal city government meeting, the leader of opposition summarily dismissed the approach:

``Computer-neeti se jan-neeti nahin banti. Is sarkar nein janta aur neta ke beech computer la diya hai (computers cannot be policy makers. This government has brought a computer between citizens and representatives).''\footnote{Statement by Ramdyal Prajapati, 18 August 2013. Accessed from the Archives of the Bhopal Nagar Nigam Record Room}

This sentiment reflected a deep-seated skepticism towards the digitalization of governance, perceiving it as a threat to the traditional approach to local politics. Councilors, in particular, protested vehemently against what they saw as an erosion of their powers and the centralization of services. Even the Mayor, despite being affiliated with the same party as the state government, voiced concerns. This opposition from the mayor highlighted the tension between local political interests and state-level initiatives, demonstrating that party alignment does not necessarily guarantee support for centralized reforms. While the city had theoretically implemented MIS, paper registers in ward offices continued to be used for logging complaints, and attendance was still recorded through manual thumb impressions. This reversion to old methods effectively rendered the entire MIS investment redundant, as local networks successfully resisted the imposed changes.

When the new Commissioner took charge of the Corporation, the sloppy implementation of the MIS came into question. As was the practice with JNNURM, Delhi threatened to hold back potential funding, were the city to not go ahead and implement change. Commissioner Das's understanding from the time was:

``You can understand the implications of the MIS being jeopardized. We would have lost our money. Negotiating is not my job, but I did ask my deputies to look into this – they knew the landscape better. The guidelines were clear; get the MIS running.''\footnote{Interview with Priyanka Das, Indian Administrative Services, Madhya Pradesh Cadre. Bhopal 2022}

The task of operationalizing the MIS fell to the Project Implementation Unit (PIU), led by State bureaucrats\footnote{Interview with Neelesh Dubey, Madhya Pradesh Municipal Administrative Services. Bhopal 2022} The Project Manager of the PIU, along with their colleagues (other state level officers) were asked to engage with protesting councilors and the Mayor to find a middle ground.\footnote{ibid. These versions were triangulated with consultants working in the PIU. Interview with Sanjay Saxena, Managing Director at Total Synergy Consulting, Delhi 2022} The core issues revolved around the redistribution of power: councilors were reluctant to relinquish control over manpower allocation and the ability to deploy sanitation workers at their discretion, while the MIS aimed to centralize these functions.

The state level bureaucracy made use of its networks existing with the councilors to convince them of the benefits of the MIS. These connections had been fostered over years of working with them at the zone-level meetings, where once every two months, the councilors of twelve wards would sit with the state level officers and lower bureaucracy in public to discuss issues of their areas. These zonal meetings brought lower bureaucracy, councilors and the state level officers in the same space, providing an opportunity for the bureaucrats to highlight how the problems could be solved if they were using the MIS.

Strategies to convince politicians included showing them how the MIS was playing out in Mumbai and Ahmedabad, convincing them that a phased implementation of the MIS would allow for a gradual transition from the old system to the new one, and suggesting a hybrid model that would preserve some level of local decision-making, particularly in the deployment of sanitation workers, while still integrating with the centralized MIS.\footnote{Interview with Bruce Pollock, who headed the DFID program in Madhya Pradesh from 2008 to 2013. Zoom 2022, 2023} The PIU team also leveraged the support of the Department for International Development (DFID) to fund the expansion of ward offices across the city.\footnote{ibid.}

When the negotiations hit a stalemate, an intervention from higher authorities, including the Chief Minister's office was sought. The State bureaucrats leveraged the presence of their co-cadre officer in key positions, particularly, that of the Officer on Special Duty (OSD) to the Chief Minister, to exert pressure on the local politicians.\footnote{Interview with Abhay Ranjan Gowankar, Madhya Pradesh Provincial Service Commission. Bhopal 2022} The intervention of the Chief Minister's office (CMO) was key to breaking this deadlock. Explicitly, one can argue that it was the fear of sanction of their political bosses which brought the dissenting Mayor and councilors to agree, but the point here is about the sustained and persistent efforts of the state bureaucracy without which, this sanctioning effect would not have been possible.\footnote{It is important to also note that, unlike in Lucknow where the state bureaucrats reached out to politicians to intervene, it was the intra-bureaucratic coordination which was key in this case.}

The 'agreement to allow' implementation of the MIS in a phased way, allowed councilors to retain some discretion in deploying sanitation workers, preserving a degree of their traditional authority. In exchange, all employees, including these workers, were required to participate in the MIS reporting and attendance system. This compromise represented a delicate balance between maintaining local political dynamics and advancing administrative gains. Furthermore, the agreement included the expansion of ward offices to all wards in the city, a significant increase from the previous 16.\footnote{Interview with Devinder Chauhan, Madhya Pradesh State Administrative Services, Bhopal, 2022. This version was corroborated by Bruce Pollock and Jim Arthur, both of whom were DFID officers in Madhya Pradesh. Zoom 2022, 2023.} This expansion, funded by the Department for International Development (DFID), served as a sweetener for the councilors, providing them with enhanced local infrastructure.\footnote{DFID invested in the creation of a new ward. Madhya Pradesh Urban Services for the Poor (MPUSP) Project. (2012). Project completion review: Madhya Pradesh Urban Services for the Poor (MPUSP) Project. Retrieved from https://iati.fcdo.gov.uk/iati\_documents/3780260.odt. Also, as the councilor from Misrod ward told me in Bhopal, 2022, ``I accepted because the office would be manned by my people.''} The land allocation for these new offices was facilitated through negotiations with the Chief Minister's office, a process expedited by the presence of State bureaucracy in key positions, including as Officer on Special Duty (OSD) to the Chief Minister.\footnote{Interview with Devinder Chauhan, Madhya Pradesh State Administrative Services, Bhopal, 2022.}

Dubey's reflection on this process is key:

``Dekho, MIS implementation ke samay teen Commissioner aaye. Gaur Mam, Gharpale Sir aur Manish Singh Sir. Teeno alag department se aaye aur alag department gaye.
Agar hum yaha iska implementation nahin dekhte to sarkar mein to wahi sawal bar bar aate hai: file kaha hai, kya matter hai, meetings hoti, kuch nahin nikalta. Hamare hone se hume pech kaha phas raha tha, yeh to pata tha. Zameen shasan ne di, computer or kendra kholne ka paisa DFID se le liya gaya.'' (Look, during the MIS implementation, three Commissioners came: Gaur Ma'am, Gharpale Sir, and Manish Singh Sir. All three came from different departments and went to different departments. If we hadn't overseen the implementation here, the same questions would keep coming up: where is the file, what's the matter, meetings would happen, nothing would come out. Because we were here, we knew where the problem was stuck. The government gave the land, and the money for computers and opening centers was taken from DFID.)\footnote{Interview with Neelesh Dubey, Madhya Pradesh Municipal Administrative Services. Bhopal 2022}

The successful implementation of the Management Information System (MIS) in Bhopal was the result of a four-year negotiation process spanning three different Commissioners. The State bureaucrats played a crucial role in this process, leveraging their unique position and capabilities in ways that central bureaucrats, bound by rule-based law and limited negotiation powers, could not. These state officials, through their long-term presence in the system, accumulated valuable institutional memory, local knowledge, and relationships. Acting as ``boundary spanners'' \citep{Tushman1981}, they operated within the state apparatus while negotiating with the political class and utilizing influence networks. Their latent information proved critical in navigating the complex political landscape, building relationships with councilors, and securing necessary resources and support from the government and external agencies like DFID. Unlike central bureaucrats limited by an expected apolitical stance and ``unwritten'' boundaries, the State bureaucracy was ideally positioned. They had fostered complex relationships with councilors through zone meeting interactions. This dynamic allowed for open communication and access between officers and councilors, while also maintaining the ability to impose sanctions if necessary. This dual nature of their relationship enabled effective negotiation and implementation of the MIS project. The state bureaucrats' awareness of the state government's (shasan) commitment to expanding ward offices and available international financial institution (IFI) budgets for infrastructure modernization further contributed to their effectiveness in aligning various departments and completing complex, multi-faceted delivery programs.

Secondly, the state bureaucracy had an interstitial role that allowed them to move beyond their formal responsibilities. They were trusted by both the senior bureaucracy and the politicians to negotiate and find a solution. When I asked the Commissioner who said ``they [state bureaucracy] knew the landscape better'', how would she know if the state level bureaucrat would remain within the mandate given and not cut deals with the politician, they replied:

``I write their ACRs are a crucial monitoring tool for the Indian bureaucracy, serving as the primary method for annual performance review of civil servants.\footnote{Purohit, B., \& Martineau, T. (2016). Is the Annual Confidential Report system effective? A study of the government appraisal system in Gujarat, India. Human Resources for Health, 14(1), 33. https://doi.org/10.1186/s12960-016-0133-8} The process involves three key components: authorship by a senior officer, review by the subject officer, and approval by a higher authority. While designed to assess performance, facilitate career advancement, and aid in human resource development, the ACR system is notable for its confidentiality and the significant impact of negative entries. The process typically includes a self-appraisal, an overall grade, and the opportunity for the appraisee to contest adverse remarks. Despite its uniform format across various roles, the ACR's importance lies not just in its content, but in the norms and perceptions surrounding it within the bureaucratic structure.

Taken together, this section highlights the crucial boundary-spanning role of state elite bureaucrats in implementing complex reforms like the MIS in Bhopal. These bureaucrats effectively bridged high-level policy directives with ground-level realities, leveraging their deep local knowledge and relationships to navigate the complex political landscape. Their unique position allowed them to synthesize disparate elements—past commitments, available resources, and political interests—into a cohesive, implementable plan. Simultaneously, the autonomy and sanctioning powers of the central bureaucracy, exemplified through mechanisms like Annual Confidential Reports (ACRs), provided a necessary system of checks and balances. This interplay between the flexible, adaptive approach of state bureaucrats and the oversight of central authorities created a framework that enabled successful reform implementation while maintaining accountability. Ultimately, this case underscores the importance of recognizing and utilizing the complementary strengths of different bureaucratic levels in the functioning of the state.

\subsection{MIS in the field}

The implementation of the Management Information System (MIS) revealed complex on-the-ground realities. This case study examines the disparities in complaint resolution across socioeconomic groups, the interplay of formal and informal channels, and the empowerment of state level officers offering insights into how the reform played out.

In Bhopal's Indira Gandhi ward, a small informal colony, lies adjacent to India's first private air-conditioned railway station. Locally known as ``Railway patri ke par'' (opposite the railway track), the tarpaulins and walls of houses here were covered in dust from Kota stone being chiseled for the temple under construction at the slum's boundary.

Here, in October 2021, I met Sartaj who struck up a conversation watching me click pictures of the BMC's sewage cleaning suction truck, which the basti owners had called to clean up the blocked-up sewer line. Over subsequent months, I observed the area to assess the city government's claim of a trash free city. The slum's proximity to a commercial hub ensured regular trash collection. However, this was supplemented by informal waste segregation within the basti, where residents sorted valuable recyclables, leaving the remaining waste in mounds outside. Despite causing occasional drain blockages and odor issues, especially in summer, sanitation workers typically collected this waste every three days.

During the Navratri festival, when the city is lit up, the entry to basti was clogged with trash which had not been picked up for a week now. Ramesh, who worked in the tea stall close by, explained to me the informal process of complaint resolution. Residents of the basti would reach out to their local neta (leader) like Sartaj, who would then contact the ward officer. If the matter was not solved, the matter would escalate to the jhuggi prokosht adhyaksh (Slum unit incharge), who would then approach the councilor. The councilor's typical approach to addressing issues involved first contacting lower-level officers, specifically the zone in-charges. If the situation escalated, they would then reach out to state level bureaucrats. However, in this instance, the councilor was away attending a sammelan (conference) in Varanasi. In their absence, Sartaj turned to the newly implemented city government application. This digital platform allowed users to report problems by capturing the location, time, and nature of the issue, automatically uploading this information to the city's website for processing. Next day, Sartaj's phone showed that the complaint had been solved and was hence closed. As the matter lingered, visits to the ward office were made where they were told that owing to the festive week, they were short-staffed, and the issue would get promptly addressed in the coming days. Meanwhile, the basti employed private people from within to collect the trash and throw it on the railway track at night, only to be promptly shooed away by the Railway police officers. The situation deteriorated, Sartaj received a perfunctory call center follow-up, where he was met with a standard reassurance: ``Nishchint rahiye, apki samasya ka samadhan jald hoga. Har mahine ke chauthe guruwar ko zone nirikshan hota hai. Aap waha bhi nagar nigam ke samne apni samasya rakh sakte hai.'' (Rest assured, your problem will be solved soon. Every fourth Thursday of the month is when the Zone in-charges visit the ward for inspection. You can also present your problem to the municipal corporation representatives there.)

At the zone meeting held in the Indira Gandhi ward, I observed a significant gathering of residents from the basti and the business in the area. Both the Zone In-charge and the Assistant Commissioner of the municipal government were present. As the Zone In-charge reviewed the area's ongoing developmental work, the local BJP councilors' right hand man in the area intervened, ``Sir, yeh area Bhopal ka main area hai, aur abhi Navratre ke time market ke samne kuda bhara pada hai'' (Sir, this is Bhopal's main area, and currently, during Navratri, there's accumulated waste in front of the market).

The sanitation officer, Ajay Sindhi, responded by asserting that routine cleaning was happening and cited staff shortages as a current challenge. Sartaj promptly interjected, ``hamare area ka bhi yehi haal hai'' (Our area is in a similar condition), and quickly flipped his phone out to show pictures to those present. He added ``aur sir, app pe complaint bhi closed dikha raha hai'' (the complaint status is shown as closed in the app). The Assistant Commissioner looked at Sindhi, questioning, ``farzi report kyun band kar rahe ho?'' (Why are you filing a closed report for unresolved complaints?). Sindhi attempted to explain further, citing staff shortages due to the festive season and assuring prompt cleaning for the upcoming festive week.The Assistant Commissioner after consulting a file provided by his assistant, stated, ``pure sheher mein apke area se sabse jyada complaint hai aur MIS mein aayi yaha 85-90\% samsyaein closed kar de gay hai'' (Your area has generated the highest number of complaints in the entire city and yet, Management Information System (MIS), shows 85-90\% complains have been marked as resolved.) As Chauhan, tried to explain that it is a ``commercial area'' and that he needs more people, the Commissioner stopped and said ``abhi pehle apke khilaf karyawahi hogi, fir age ki karywahi ki jayegi'' (First, action will be taken against you, then further steps will be taken).

Sindhi was still seen at the ward office the following week, having been issued only a show cause notice, with a suspension to follow a couple of weeks later after the investigation concluded.\footnote{Raj Express. (n.d.). Bhopal Municipal Corporation health officer suspended for taking bribe. https://www.rajexpress.com/india/madhya-pradesh/bhopal-municipal-corporation-health-officer-suspended-for-taking-bribe} Interestingly, while the issues in Ramesh's area persisted, there was a noticeable change in the responsiveness of local staff. They became more alert to complaints coming from Sartaj or the area, ensuring that weekly complaint registers were duly collected and responded to promptly.\footnote{The registered complaint from the ward saw a reduction by 40\% in the subsequent weeks when Sindhi had left. Data collected for the period of September, October and November 2022 from the Bhopal Nagar Nigam Zone 6 office.}

The implementation of the Management Information System (MIS) in Bhopal's trash management demonstrates the potential of technological solutions only when integrated with robust audit practices. To compare, in Lucknow where concessionaires operated with minimal oversight mechanisms, the absence of such integrated approaches led to suboptimal outcomes. This contrast underscores the importance of combining technological interventions like MIS with close monitoring and responsive feedback systems to ensure effective urban governance.

Further, the fieldwork in four wards revealed divergent patterns in complaint resolution. MP Nagar, predominantly comprising planned colonies and inhabited by a socioeconomically advantaged population, demonstrated effective utilization of formal complaint mechanisms, such as mobile applications and direct bureaucratic connections. In contrast, Bawadiya Kalan, mostly comprising migrant populations, relied heavily on local councilors and intermediaries for channeling grievances. This disparity underscored a critical challenge in these technocratic reforms: integrating complaints from diverse social groups across both formal and informal networks.

The discrepancies between MIS reports and ground realities, as observed during the inspection of Indira Gandhi Ward, underscored the crucial role of state bureaucracy in addressing implementation gaps. The incident involving Sindhi and the Assistant Commissioner illustrated a significant shift in bureaucratic dynamics, where accountability took precedence. The Assistant Commissioner's decisive action in confronting Sindhi over false reporting and initiating immediate disciplinary measures had a ripple effect, signaling the importance of accurate reporting and accountability within the system. As Sangeeta Tiwari, a lower division clerk in the same zone, aptly remarked,

``MIS mein number idhar-udhar ho sakte hain, par jab log khud yahan aa jate hain, tab kya jawab denge? Isliye hume dhyan rakhna padta hai ki MIS sirf ek auzar hai, asli jawabdehi to zameen par hota hai'' (Numbers in MIS can be manipulated, but when people show up here, what answer will we give? That's why we must remember that MIS is just a tool; real accountability happens on the ground)

\section{``Stickiness'' of State Capacity in Bhopal}

How can we determine if there was an incremental addition to state capacity in Bhopal? A true measure of incremental state capacity would lie in the city government's ability to continue building capacity and address more complex challenges. Bhopal's ability from basic trash collection to tackling more intricate issues such as legacy waste management and green climate bond monetization are good examples of building of this incremental state capacity. I discuss the issue of legacy waste here.

Bhanpur Khanti is a 36.90-acre area, used exclusively for waste dumping since 1970, and had accumulated over 20,000 tons of waste by 2018.\footnote{Down To Earth. (2022). Cleanest cities of India: Bhopal reclaimed 37 acres of wasteland by clearing legacy waste. Retrieved from https://www.downtoearth.org.in/waste/cleanest-cities-of-india-bhopal-reclaimed-37-acres-of-wasteland-by-clearing-legacy-waste-81654} In 2021, during my visit, the last remnants of legacy waste was being processed and by 2023 summer, the site had been cleared and replaced with a garden. Residents of the field site told me that the groundwater, though still an issue, had improved.\footnote{Interview with Ajay Pannu, resident of Khejrabaramad area in the Bhanpur ward of the Bhopal Municipal Corporation (BMC). Bhopal 2023}

The project took about two decades to complete. First, in a resource constrained city and state, cleaning legacy waste was not a priority.\footnote{Interview with Vivek Aggarwal, Swachh Bharat Mission Incharge, Urban Development and Housing Department, Madhya Pradesh Government, Bhopal, 2023.} Second, the processing of this waste required the employment of sensitive technology - the dumpsite had various kinds of waste, each of which required segregation and then processing. The only accessible funding was from the union government.

The proposal originated from the Commissioner's office back in 2014, but by the time the proposal went through various levels of government, a new Commissioner had joined.\footnote{Interview with Vivek Aggarwal, Swachh Bharat Mission Incharge, Urban Development and Housing Department, Madhya Pradesh Government, Bhopal, 2023.} In 2016, when the project was approved and tenders floated, no firm showed interest.\footnote{Reply filed by the Bhopal Nagar Nigam to the Central pollution Control Board. National Green Tribunal. (2021). Compliance of Municipal Solid Waste Management Rules, 2016 and other environmental issues (Original Application No. 606/2018). Central Pollution Control Board. Retrieved from [https://cpcb.nic.in/uploads/MSW/Reports\_swm\_8.pdf]} Expensive technology and long term commitment meant that the very few companies in India which handled waste at such a large proportion thought the money offered was not commensurate with the amount of trash. The state bureaucrats worked with these firms, understood the problem and brought this problem to light in meetings with the central bureaucracy.\footnote{Interview with Vivek Aggarwal, Swachh Bharat Mission Incharge, Urban Development and Housing Department, Madhya Pradesh Government, Bhopal, 2023.} Finally, Saurashtra Environment was awarded the bid in 2018. In the middle of the work in 2020, the company increased the amount for waste segregation from the tender amount, stating that they had found traces of hazardous chemicals originating from the 1984 Union Carbide chemical disaster.\footnote{See report filed with the CPCB:https://cpcb.nic.in/uploads/MSW/Reports\_swm\_8.pdf} The BMC faced a shortage of money and with a new Commissioner, the same wheel had to be reinvented.

``I have hundreds of files on my table. I just cannot keep track of what happened during the tenure of the Commissioners before me. For this project, I had to completely rely on my junior officers to update me.''\footnote{Interview with Vijay Datta, Swachh Bharat Mission Incharge, Urban Development and Housing Department, Madhya Pradesh Government, Bhopal, 2023.}

The specificity and uniqueness of the project required urgent attention. Given the ``cluttered'' engagements of the city chief, it was the State bureaucracy that prioritized the request for seeking additional money for fear of Saurashtra Environment to abandon the project midway.

``First line of leadership changes constantly. But the second line of command is what is important - that is what makes the difference.''\footnote{Interview with Devendra Chauhan, Madhya Pradesh Municipal Administrative Services, who was then with the BMC and is now an Officer on Special Duty (OSD) to the Minister of Urban Development. Bhopal 2023}

The additional grant was made available and during my visit in the summer of 2023, the once mountainous trash site had been cleared and replaced with a garden. Between 2014 and 2023, during the life of Bhanpur Khanti project, Bhopal saw nine Commissioners. This rapid turnover at the top could have derailed the project, but the institutional memory within the state bureaucracy kept it on track. Acknowledging their role, a private consultant with experience working with various IFIs and city government summarized:

``Devinder Chauhan types [state level bureaucrat] and others like him would become super useless in front of the IAS, but their usefulness only comes to light when the IAS has left the room. They knew the inside out of how programs are structured and how we were doing it. If you get them on board, the project is successful; otherwise, it is not.''\footnote{Interview with Arkaja Singh, a private consultant with experience working with various International Financial Institutions and was a consultant hired by DFID for UDAY. Delhi 2022}

The case of Bhanpur Khanti underscores the pivotal role of state bureaucrats in bolstering and sustaining state capacity through their deep institutional knowledge and continuity within the bureaucratic system. This project, spanning nearly two decades (ten since the proposal was finally sent, and ten when the first ``file noting was done'', illustrates how these bureaucrats, often overshadowed by the more visible first-line leadership, were instrumental in navigating the complexities of governance that arise due to high turnover at the top. Despite the rapid succession of central level bureaucrats – the Commissioners in Bhopal, the project remained on track largely because of the institutional memory and expertise retained by these bureaucrats. Their intimate understanding of the bureaucratic processes, from handling files to knowing the exact locations of key documents and the intricacies of budget allocations, was critical in pushing the project forward. As the case highlights, their ability to maintain continuity in the face of administrative change, resolve bottlenecks, and secure necessary resources, demonstrates how their institutional knowledge directly contributes to the effective implementation of policies.

\section{Theoretical Framework: Bureaucratic Capacity in Diverse Societies}

In comparative case study research, it is crucial to conduct a thorough, context-specific analysis of each case before engaging in comparative analysis. Scholars such as \citet{Yin2017} and \citet{Gerring2007} argue that in-depth examination of individual cases allows researchers to uncover the complex, context-dependent mechanisms that drive outcomes, which might be obscured in direct comparisons. This approach ensures that each case is fully understood in its own right, providing a solid foundation for meaningful comparison.

Applying this method, this chapter first focuses on Bhopal to explore in detail the unique factors that contributed to its success in enhancing bureaucratic capacity and delivering public goods. By analyzing the mechanisms of horizontal and vertical coordination, the role of internal intermediaries, and the leveraging of institutional knowledge within Bhopal's governance framework, this section aims to uncover the intricate dynamics that underpinned the city's governance achievements.

Once this in-depth analysis is established, the study will then broaden its scope to include a comparative perspective with Lucknow. This structured approach ensures that the distinctive strategies and challenges in Bhopal are fully articulated, setting a robust foundation for understanding the divergences observed in Lucknow . By initially isolating Bhopal for detailed analysis, the study follows best practices in case study research, ensuring that the nuances of each city's governance are thoroughly explored before drawing comparative conclusions.

\subsection{State Capacity and Bureaucratic Coordination in Rapidly Urbanizing Societies}

The effectiveness of public service delivery in rapidly urbanizing societies is a central concern in political science, with scholars often debating the role of state capacity and institutional strength in addressing the challenges that arise in such contexts \citep{Alesina1999, Habyarimana2007}. While issues like social fragmentation are frequently discussed as potential obstacles, this chapter argues that the critical factor influencing outcomes is the configuration and capacity of state institutions.

Building on existing scholarship, this chapter posits that successful public goods provision in urbanizing societies is not merely a matter of overcoming social cohesion or fragmentation. Rather, it depends on the ability of state institutions to manage these challenges through robust institutional frameworks and effective state-society relations \citep{Heller2001}. The goals and actions of state officials, shaped by historical and political contexts, are pivotal in determining how state institutions navigate the complexities of urbanization \citep{Singh2015}.

Moreover, this chapter contends that a key determinant of service delivery in these rapidly changing urban environments is the capacity of the bureaucracy to coordinate and implement policies efficiently. Effective state capacity requires bureaucracies to achieve both horizontal coordination across government departments and vertical coordination between different levels of government. By examining the case of Bhopal, this chapter demonstrates how state bureaucrats enhanced bureaucratic capacity through these forms of coordination, enabling the city to deliver public goods effectively despite the evolving challenges associated with rapid urbanization.

\subsubsection{Horizontal and Vertical Coordination as Mechanisms of Bureaucratic Capacity}

The effective delivery of public services in Bhopal can be attributed to the strategic capacity of state bureaucrats, who adeptly navigated both horizontal and vertical coordination within the institutional framework. Horizontally, these bureaucrats facilitated interdepartmental collaboration, enabling communication and cooperation across state boundaries. This approach was crucial in collecting valuable feedback from frontline sanitation workers, which informed the development of standardized hiring and transfer protocols—an example of bottom-up horizontal coordination. Simultaneously, they secured a significant degree of autonomy from central bureaucratic oversight, allowing them to implement a Management Information System (MIS) with minimal external interference, demonstrating top-down horizontal coordination.

Vertically, the bureaucrats effectively bridged intra- and inter-institutional divides to ensure policy coherence and consistency. Within the institutional hierarchy, they coordinated closely with colleagues to develop and enforce uniform hiring and transfer regulations across the state, reflecting strong vertical in-institution coordination. Additionally, they skillfully engaged with external political actors, particularly local politicians, to secure the necessary support for the MIS initiative. This vertical out-institution coordination was critical for the successful adoption and implementation of the system, illustrating the bureaucrats' ability to integrate administrative processes with political considerations.

The capacity of state level bureaucrats to navigate both horizontal and vertical coordination highlights their strategic role in the effective delivery of public services in Bhopal. To understand the mechanisms behind this success, it is essential to explore the specific functions these bureaucrats perform—particularly as internal intermediaries, boundary spanners, and custodians of institutional knowledge—which collectively underpin their ability to maintain coherence, foster collaboration, and ensure continuity across the state's complex administrative landscape.

\subsubsection{First, Internal Intermediaries: Bridging Policy and Implementation}

State-level bureaucrats serve as critical internal intermediaries within the institutional framework, occupying a unique position that straddles the space between high-level policymakers and frontline service providers. Traditionally, intermediaries are conceptualized as entities that mediate between the state and its citizens, often operating outside the formal boundaries of state structures.

However, by reversing this gaze, we can see that within the state itself, these bureaucrats function as internal intermediaries. They play a pivotal role in ensuring that policy directives from the upper echelons of government are effectively translated into actionable plans at the local level. By mediating between different layers of the bureaucracy, they maintain the coherence and continuity of public service delivery, which is essential for the smooth functioning of the state. In Bhopal, this intermediary role was decisive in bridging the gap between the state's strategic goals and on-the-ground implementation, a factor that significantly contributed to the city's success in delivering public goods. These bureaucrats enhance the internal flow of information and align the actions of various governmental levels, ensuring that policy initiatives are consistently carried out even in the face of complex administrative challenges.

\subsubsection{Second - Boundary Spanners}

The success of state programs in Bhopal often hinged on the bureaucrats' ability to function as boundary spanners, facilitating coordination across different departments and engaging with external political actors. These bureaucrats excelled in transcending the formal boundaries of their roles to build relationships that were critical for aligning various stakeholders towards common goals. By engaging actively with local politicians and fostering deep relationships with councilors through regular interactions, they created an environment of mutual respect and cooperation, which was essential for the effective implementation of complex, multi-faceted delivery programs. In contrast to central bureaucrats, who are often constrained by rigid, rule-based structures, these state-level officials had the flexibility to negotiate and build informal networks of influence, which allowed them to span boundaries within the state apparatus effectively. Their ability to ``close the loop'' on policy implementation by engaging both horizontally across departments and vertically with political actors played a significant role in the successful execution of state initiatives in Bhopal, ensuring that policies were not only adopted but also adapted to local contexts.

\subsubsection{Third - Institutional Knowledge}

A deep well of institutional knowledge underpins the effectiveness of state-level bureaucrats in maintaining and enhancing state capacity. Drawing from decades of experience and a deep understanding of the bureaucratic system, these officials know exactly where critical information is stored, who controls it, and how to navigate the often-labyrinthine structures of the state. This expertise allows them to resolve bottlenecks, trace vital files, and ensure that resources are allocated efficiently, which is critical for the continuous functioning of the state. In Bhopal, this institutional knowledge was key to navigating the complexities of the bureaucratic system, enabling the state to maintain continuity in public service delivery despite challenges. The bureaucrats' ability to leverage this knowledge to drive the state's machinery forward demonstrates how institutional memory is not merely about maintaining continuity but is also about ensuring that the state remains responsive and effective. Their intimate familiarity with the intricacies of budget allocations and procedural processes further allowed them to mobilize resources efficiently, ensuring that policies were implemented effectively and that public services were delivered consistently across the state.

This study has analyzed the mechanisms through which state-level bureaucrats in Bhopal enhanced bureaucratic capacity, focusing on the critical roles of horizontal and vertical coordination, internal intermediaries, boundary spanning, and institutional knowledge. The findings emphasize how bureaucratic agency and strategic coordination are pivotal in navigating the complexities of governance in rapidly urbanizing societies.

\section{Comparisons across Cases}

This section examines the underlying political and institutional dynamics that contribute to these differences, focusing on how bureaucratic structures, political representation, and the strategic use of privatization influence the capacity of the state to manage societal demands. By comparing these two cases, the analysis sheds light on the broader implications for state capacity in diverse urbanizing contexts and offers insights into the conditions under which state institutions can effectively mediate between societal pressures and administrative imperatives.

\subsection{Political Structure and State-Society Relations: Bhopal vs. Lucknow}

Bhopal and Lucknow share a common political structure, characterized by directly elected Mayors who possess limited executive authority. In both cities, councilors, even those participating in key committees like the Karyakarini Samiti in Lucknow and the MIC in Bhopal, wield minimal power. This structural limitation forces politicians into the role of ``firefighters,'' addressing immediate concerns rather than engaging in substantive legislative or policy-making functions. The constrained functional, financial, and legislative powers of Indian cities, where trash management consumes the majority of manpower and budgets, leave little room for broader governance initiatives. As a result, political legitimacy in these urban centers often hinges on the ability to selectively distribute resources and services, reinforcing clientelist practices rather than fostering institutional capacity or public accountability.

\subsection{Enhanced Representation and State Capacity in Bhopal}

In contrast to Lucknow, Bhopal has experienced a modest yet significant enhancement in city government authority, which has had implications for state capacity and state-society relations. The presence of a functioning council that meets quarterly and the establishment of ward committees at the zonal level have introduced elements of deliberative democracy, allowing for a more inclusive decision-making process. Crucially, Bhopal provides better representation for marginalized groups, with Muslims and Dalits enjoying proportional or even over-representation at various levels of governance. This inclusivity fosters a sense of visibility and security among these communities, reducing the likelihood of identity-based politics dominating the governance agenda. As a result, politicians in Bhopal are less inclined to manipulate identities for political gain, which in turn allows bureaucrats to function more autonomously and effectively address city-level challenges.

\subsection{Fragmented Representation and Bureaucratic Constraints in Lucknow}

Conversely, in Lucknow, the representation of marginalized groups such as Muslims, OBCs, and SCs is largely confined to the city level, making it the primary platform for these communities to voice their concerns. Political parties in Lucknow have cultivated strong support bases among these groups, each emphasizing distinct political identities. However, the absence of formal mechanisms like ward committees and the irregularity of corporation meetings result in councilors becoming the sole mediators of these identities, often leading to fragmented and contentious state-society relations. This dynamic compels politicians to prioritize immediate service delivery to their constituencies, frequently pressuring bureaucrats or higher-level officials to comply with these demands. The outcome is a bureaucracy in Lucknow that is both compliant and constrained, leading to selective service distribution and a persistent inability to develop comprehensive, integrated solutions for the city's public service needs. This selective representation paradoxically undermines public service delivery, despite the presence of formal political inclusion.

\subsection{Bureaucratic Role and State Capacity: Contrasting Dynamics}

The most significant divergence between Bhopal and Lucknow lies in the role and capabilities of their respective bureaucracies. Although both cities are situated in predominantly rural states, where municipal management constitutes a small part of bureaucratic responsibilities, Bhopal's state-level bureaucracy benefits from clear service rules governing promotions and transfers. These rules create a unique dynamic where state-level officers are held accountable by IAS officers while also enjoying a predictable career trajectory. This dual mechanism of accountability and predictability yields two critical outcomes: first, it allows IAS officers to grant autonomy to state-level bureaucrats, confident in their ability to sanction any deviation from responsibilities; second, it incentivizes state-level bureaucrats to align interests across line departments, knowing that career advancement depends on their performance and coordination skills. Consequently, Bhopal's municipal cadre has been instrumental in coalition-building, navigating complex political landscapes, and acting as boundary spanners, leveraging their institutional knowledge to tackle urban challenges such as trash management. This ``second line of command,'' characterized by clear incentives and accountability, has been pivotal in Bhopal's success in enhancing state capacity at the municipal level.

\subsection{Privatization of Services and State Capacity: Diverging Outcomes}

A critical aspect of the comparative analysis between Bhopal and Lucknow is the role of privatization in shaping state capacity. The recent focus on urban development in India, driven by federally supported programs, has introduced a new era of ``Consultancy Urbanism,'' where international and domestic firms provide private expertise to implement policies at various levels. However, the experiences of Bhopal and Lucknow illustrate that the impact of privatization on state capacity is contingent on how it is integrated into local governance structures. In Bhopal, international financial institutions like the Asian Development Bank (ADB) and the Department for International Development (DFID) committed to long-term investments aimed at both infrastructure development and capacity building within the city government. This strategic engagement facilitated the alignment of interests among various actors, enabling the private sector to complement the state's efforts in areas ranging from everyday tasks like waste collection to more complex initiatives such as issuing green bonds. The success of Bhopal's privatization efforts highlights the importance of using private capital to enhance, rather than circumvent, state capacity.

In contrast, Lucknow's approach to privatization has been largely counterproductive. Instead of utilizing private sector resources to build in-house capacities, Lucknow's administration treated privatization as a shortcut to bypass the more challenging task of internal reform. This misalignment between the state's objectives and the private sector's role resulted in a dysfunctional system where bureaucrats, fixated on meeting the performance metrics of federally funded programs, neglected the deeper reforms needed to strengthen the local state. At the same time, the private sector had to navigate a fragmented labor force dominated by the local politician-bureaucrat nexus, leading to a system where negotiation and compromise, rather than efficiency and reform, became the norm. This ``endpoint solution'' approach, where privatization was seen as an easy fix, ultimately masked underlying systemic issues and failed to address the root causes of urban challenges. The contrasting experiences of Bhopal and Lucknow underscore the critical role of state commitment to building long-term capacities and establishing robust regulatory frameworks as prerequisites for successful privatization.


\section*{Annexure -1 What does the Indian city do?}

\begin{table}[ht]
\centering
\begin{tabular}{|l|c|c|c|c|c|c|c|c|c|}
\hline
\textbf{States} & \textbf{Urban} & \textbf{Land-} & \textbf{Roads and} & \textbf{Water} & \textbf{Solid waste} & \textbf{Slum} & \textbf{Parks \&} & \textbf{Burials and} & \textbf{Birth and} \\
 & \textbf{Planning} & \textbf{use} & \textbf{bridges} & & \textbf{management} & \textbf{Improvements} & \textbf{Gardens} & \textbf{burial grounds} & \textbf{Death Registration} \\
\hline
Andhra Pradesh & M & M & M & M & C & M & C & C & C \\
\hline
Bihar & M & M & M & C & C & C & M & C & C \\
\hline
Delhi & M & S & M & S & C & M & M & C & C \\
\hline
Gujarat & S & S & M & M & C & M & M & C & C \\
\hline
Haryana & S & S & M & M & C & M & C & C & C \\
\hline
Karnataka & M & M & M & M & M & M & C & C & M \\
\hline
Kerala & M & M & M & S & C & M & M & C & C \\
\hline
Madhya Pradesh & M & M & M & M & C & C & C & C & C \\
\hline
Maharashtra & M & M & M & C & C & M & C & C & C \\
\hline
Odisha & M & M & M & M & C & C & M & C & C \\
\hline
Punjab & S & S & M & M & M & M & M & C & C \\
\hline
Rajasthan & M & S & M & S & C & M & M & C & C \\
\hline
Tamil Nadu & M & M & M & M & C & M & M & C & C \\
\hline
Telangana & M & M & M & M & C & M & M & C & C \\
\hline
Uttar Pradesh & S & S & S & M & C & S & M & M & M \\
\hline
West Bengal & M & M & M & C & C & M & M & C & M \\
\hline
\end{tabular}
\end{table}

[Annual Confidential Reports]. If they are mis-aligned I have the power to reproach''\footnote{Interview with Priyanka Das, Indian Administrative Services, Madhya Pradesh Cadre. Bhopal 2022}

ACRs

